\documentclass[12pt,oneside]{book}
\usepackage{graphicx}
\usepackage{epsfig}
\usepackage{amssymb}
\usepackage{amsmath}
%\usepackage{amsmath}
\usepackage{amsthm}
\usepackage{longtable}
\usepackage{slashbox}
\usepackage{multirow}
\usepackage{tabularx}
%\usepackage{amssymb}
\usepackage{graphicx}
\usepackage{graphics, t1enc}
\usepackage[spanish]{babel}
\usepackage{makeidx}
%\usepackage{titlesec}
\makeindex
\makeglossary

\usepackage{pstricks}
%\usepackage{showkeys}
\setlength\unitlength{1cm}
\newif\ifAMS
\IfFileExists{amssymb.sty}
{\AMStrue\usepackage{amssymb}}{}
\usepackage{latexsym}



\renewcommand{\baselinestretch}{1.5}
\textwidth=16.5cm
\textheight=21.5cm
\footskip=1.5cm
\oddsidemargin -0.6cm
\voffset =-1cm
%\renewcommand\chaptername{Cap\' {\i}tulo}
%\renewcommand\bibname{Bibliograf\' {\i}a}

\newtheorem{teo}{\sc Teorema}[section]
\newtheorem{ejemplo}{\sc Ejemplo}[section]
\newtheorem{lemma}{\sc Lema}[section]
\newtheorem{obs}{\sc Observaci\'on}[section]
\newtheorem{coro}{\sc Corolario}[section]
\newtheorem{definition}{\sc Definici\'on}[section]
\newtheorem{prop}{\sc Proposici\'on}[section]
\newtheorem{nota}{\sc Nota}[section]
\newcommand{\ese}{\mathcal{S}}

\newtheorem{notacion}{\sc Notaci\'on}[section]

\newcommand{\CC}{\mathbb{C}}
\newcommand{\NN}{\mathbb{N}}
\newcommand{\PP}{\mathbb{P}}
\newcommand{\HH}{\mathbb{H}}
\newcommand{\pp}{\mathbb{\overline{P}}}
\newcommand{\DD}{\mathbb{D}}
\newcommand{\QQ}{\mathbb{Q}}
\newcommand{\RR}{\mathbb{R}}
\newcommand{\ZZ}{\mathbb{Z}}
\newcommand{\EE}{\mathbb{E}}
\newcommand{\TT}{\mathbb{T}}
\newcommand{\XX}{\mathbb{X}}

\newcommand{\II}{\mathbb{I}}

\newcommand{\der}{\mathcal{D}}
\newcommand{\kk}{\mathcal{K}}
\newcommand{\mm}[1]{{\mathcal{M}\/}(#1)}
\newcommand{\nequiv}{{\equiv \hspace*{-3.7mm}/}}
\newcommand{\capa}[1]{\mbox{{\em Cap\/}}(#1)}
\newcommand{\gra}[1]{\mbox{{\em grad\/}}(#1)}
\newcommand{\sop}[1]{\mbox{{\em supp\/}}(#1)}
\newcommand{\esup}[1]{{\mbox{\rm ess\,sup\/}}#1}
\newcommand{\lqqd}{\hfill $\blacksquare$}

\renewcommand{\thefootnote}{\fnsymbol{footnote}}
\newcommand{\ca}[1]{{\em card\/}(#1)}
\newcommand{\dis}[1]{{\em dist\/}(#1)}
\newcommand{\imm}[1]{{\em Im\/}#1}
\newcommand{\ree}[1]{{\em Re\/}#1}
\newcommand{\conv}[1]{{\em conv\/}(#1)}
\newcommand{\uni}{\; \; \overrightarrow{\longrightarrow}\; \;}
\newcommand{\unin}{\overrightarrow{\longrightarrow}}
\newenvironment{ob}{\begin{obs}{\rm} \end{obs}}

\def\bbuildrel#1_#2^#3{\mathrel{
 \mathop{\kern 0pt#1}\limits_{#2}^{#3}}}
\def\bbbuildrel#1_#2{\mathrel{
 \mathop{\kern 0pt#1}\limits_{#2}}}
\def\limsup{\mathop{\overline{\rm lim}}}
\def\liminf{\mathop{\underline{\rm lim}}}
%%%% redise\~nar cap\'itulos
%\usepackage{titlesec}
%\newcommand{\bigrule}{\titlerule[0.3mm]}
%\titleformat{\chapter}[display]
%{\bfseries\Large}
%{
% \titlerule
% \filleft
% \large\chaptertitlename\
% \large\thechapter}
%{0mm}
%{\filleft}
%[\vspace{0.5mm} \bigrule]
%%% formato de encabezado del cap\'itulo
\usepackage[Lenny]{fncychap}
\ChNameVar{\fontsize{12}{13}\usefont{OT1}{phv}{m}{n}\selectfont}
\ChNumVar{\fontsize{50}{52}\usefont{OT1}{ptm}{m}{n}\selectfont}
\ChTitleVar{\normalsize\bfseries}
\ChNameLowerCase
 \ChRuleWidth{1pt}






%capitulos centrados%

%\usepackage{fncychap}
%\renewcommand{\thechapter}{\Roman{chapter}}
%\ChTitleUpperCase
%\ChNameVar{\centering \bfseries}
%\ChNumVar{\centering \bfseries}
%\ChTitleVar{\centering \bfseries}

%numeracion superior%
\usepackage{fancyhdr}
\fancyhf{}
\fancyhead[R]{\thepage}
\renewcommand{\headrulewidth}{0pt}

\usepackage{hyperref}%  no mover de aqu\'i
\begin{document}
%\pagenumbering{roman}%paginas en romano%
%\pagenumbering{roman}
\frontmatter
\thispagestyle{empty}


\vspace*{-0.5cm}

\begin{figure}[h]
\includegraphics[width=3.5cm, height=1.5cm]{./ps/logo3}\ \ \ \ \ \ \ \ \ \ \ \ \ \ \ \ \ \ \ \ \ \ \ \ \ \ \ \ \ \ \ \ \ \ \ \ \ \ \ \ \ \ \ \ \ \ \ \ \ \ \ \ \ \ \ \ \ \ \ \ \ \ \ \ \ \ \ \ \ \ \ \ \ \ \ \ \ \ \ \ \
\includegraphics[width=2.5cm, height=1.5cm]{./ps/escudoua}
  \end{figure}

\vspace{3cm}
\begin{center}
\begin{large}
{\bf\large OPERADORES LINEALES EN ESPACIOS CON M\'ETRICA INDEFINIDA}
\end{large}
\end{center}
\vspace{2.5cm}
\begin{center}
\begin{large}
\textbf{Willian A. Vides Ramos}\\
\end{large}
\end{center}



\vspace{2.5cm}

\begin{center}
\textbf{Trabajo de Maestr\'ia}
\end{center}

\vspace{2.5cm}



%\end{center}



%\hspace{6cm}
%\begin{minipage}[t]{8cm}




\begin{center}
 Marzo, 2014
\end{center}

\newpage 
\thispagestyle{empty}


\begin{center}

\vspace*{-1.5cm}

\begin{figure}[h]
 \centering
   \includegraphics[width=3.5cm, height=3.5cm]{escudoua}\\
  \end{figure}

\vspace{-0.8cm}

\begin{center}
\textbf{UNIVERSIDAD DEL ATL�NTICO}

\vspace{-0.1cm}

PROGRAMA DE MATEM�TICAS
\vspace{-0.1cm}

MAESTR\'IA EN MATEM\'ATICAS

\vspace{-0.1cm}

\end{center}

\vspace{2.5cm}

\begin{large}
{\bf SOBRE UNA NUEVA CLASE DE POLINOMIOS TIPO APOSTOL GENERALIZADOS Y ALGUNAS PROPIEDADES}
\end{large}


\vspace{2cm}


\begin{large}
Trabajo Especial de Grado presentado a la Universidad del Atl�ntico por\\
\textbf{Pedro Hern�ndez Llanos}\\
\vspace{1 cm}
como requisito parcial para optar al grado de\\
\textbf{Magister en Ciencias Matem\'aticas}\\

\vspace{1cm}

Con la asesor\'ia del profesor\\

%\vspace{0.2cm}

\textbf{Alejandro Urieles G.}\\
\end{large}


\vspace{2.5cm}

Marzo, 2014
\end{center}


\thispagestyle{empty}

\addcontentsline{toc}{chapter}{Aprobaci\'on  del Jurado}

\vspace*{-1.5cm}


\begin{figure}[h]
\includegraphics[width=3.5cm, height=1.5cm]{./ps/logo3}\ \ \ \ \ \ \ \ \ \ \ \ \ \ \ \ \ \ \ \ \ \ \ \ \ \ \ \ \ \ \ \ \ \ \ \ \ \ \ \ \ \ \ \ \ \ \ \ \ \ \ \ \ \ \ \ \ \ \ \ \ \ \ \ \ \ \ \ \ \ \ \ \ \ \ \ \ \ \ \ \
\includegraphics[width=1.8cm, height=1.8cm]{./ps/logo2}
  \end{figure}


\vspace{1cm}
\begin{center}

\textbf{UNIVERSIDAD DEL ATL�NTICO}

\vspace{-0.1cm}

PROGRAMA DE MATEM�TICAS
\vspace{-0.1cm}

MAESTR\'IA EN CIENCIAS MATEM\'ATICAS

\vspace{-0.1cm}
\end{center}

\vspace{0.8cm}

\begin{center}
{\bf DISE�OS \'OPTIMOS EN LA PRESENCIA DE EFECTOS DE BLOQUES ALEATORIOS}
\end{center}


\vspace{0.1cm}

\hspace{6.5cm}
\begin{minipage}[t]{8.5cm}
\begin{large}
Estudiante: EDDIE RODRIGUEZ BOSSIO\\ C\'odigo: 72231783
\end{large}
\end{minipage}

\vspace{0.2cm}



\vspace{0.5cm} 
Este Trabajo Especial de Grado ha sido aprobado en nombre de la
Universidad del Atl�ntico  por el siguiente jurado examinador:
\vspace{0.5cm}

\begin{center}
\underline{\hspace{8cm}}

\vspace{0.1cm}

\textbf{Dr. Julio Cesar Romero Pab�n
 \\Universidad del Atl�ntico - Colombia}

\vspace{0.8cm}

\underline{\hspace{8cm}}

\vspace{0.1cm}

\textbf{M.Sc. Alejandro Sanchez Salazar
 \\Universidad del Atl\'antico - Colombia}

\vspace{0.8cm}

\underline{\hspace{8cm}}

\vspace{0.1cm}

\textbf{Dr. rer. nat. Jes�s Alonso Cabrera\\ Tutor}


\vfill

\textbf{Septiembre de 2015}


\end{center}


\chapter*{\textcolor{blue}{Agradecimientos} }

\addstarredchapter{Agradecimientos}

\protect\thispagestyle{empty}


\lettrine[lines=4,loversize=-0.1,lraise=0.1,lhang=.2]{E}{}{n primer lugar doy gracias a Dios dador de la vida y oportunidades;
quien puso en mi camino a personas e instituciones que fueron piezas
claves para este logro en mi vida. Ellos son mi familia, mi tutor
de tesis Dr. rer. nat. Jes�s Alonso Cabrera por compartir sus conocimientos
y consejos, Dr. Jorge Rodriguez, Dr. Alejandro Urieles, profesores
de la maestr\'{i}a, compa�eros de estudio y las instituciones de la
Universidad del Atl�ntico donde curs� mi pregrado y maestr\'{i}a,
y la Universidad del Norte donde realic� mi especializaci�n en matem�ticas
y se llevaron a cabo las asesor\'{i}as de este trabajo de grado, ya
que mi asesor es docente de planta de esta alma mater.}

\vspace{0.2cm}
 \hspace{11cm} Eddie Rodriguez



\thispagestyle{empty}

\ \vfill%

\hfill%
\draftbox{\begin{minipage}[c]{10cm}
%\begin{flushleft}
{%
\sffamily%
\footnotesize
\vspace{2cm}
\noindent 
\chapter*{\color{blue}AGRADECIMIENTOS}
\vspace{1.0cm}
A mis padres y a Diana Patricia, Adriana y Diana Luc\'ia .
%\end{flushleft}
\end{minipage}}

\vfill\ %


\addcontentsline{toc}{chapter}{RESUMEN}

\chapter*{RESUMEN}

\vspace{1.0cm}


\hspace{0.5cm}
En este trabajo estudiamos una nueva clase de polinomios tipo Apostol generalizados $\mathcal{Q}_n^{[m-1,\alpha]}(x,b,c;\lambda;u,v)$ con $(\alpha, u, v\in \CC$ y $b, c \in \RR^{+})$ los cuales est�n definidos en un entorno adecuado de $t=0$ por la siguiente funci�n generatriz:
\begin{equation}
\displaystyle\left(\frac{(2^ut^v)^{m}}{\lambda b^t+\sum\limits_{l=0}^{m-1}\frac{(t\log b)^l}{l!}}\right)^{\alpha}c^{xt} =\displaystyle\sum\limits_{n=0}^{\infty}
\mathcal{Q}_n^{[m-1,\alpha]}(x,b,c;\lambda;u,v)\frac{t^n}{n!}; \quad |t\log b|<|log(-\lambda)|,
\end{equation}
estos polinomios generan a las nuevas clases de todos los polinomios conocidos de Apostol, as�:\\
\begin{align}
& \mathcal{B}_n^{[m-1,\alpha]}(x;b,c;\lambda)=(-1)^{\alpha}\mathcal{Q}_n^{[m-1,\alpha]}(x;b,c;-\lambda;0,1),\\
& \mathcal{E}_n^{[m-1,\alpha]}(x;b,c;\lambda)=\mathcal{Q}_n^{[m-1,\alpha]}(x;b,c;\lambda;1,0),\\
& \mathcal{G}_n^{[m-1,\alpha]}(x;b,c;\lambda)=\mathcal{Q}_n^{[m-1,\alpha]}(x;b,c;\lambda;1,1).
\end{align}

Establecemos algunas propiedades b�sicas para estos polinomios, incluyendo relaci�n de recurrencia y la ecuaci�n diferencial que estos polinomios satisfacen. Finalmente se determinan f�rmulas de conexi�n entre los polinomios y los polinomios de Genocchi, los polinomios de Jacobi, los polinomios de Hermite, los polinomios de Laguerre, los polinomios de Charlier, los polinomios de Bessel, los polinomios de Bernoulli generalizados $B_n^{[m-1]}(x)$ y los n�meros de Stirling de segunda clase, lo cual extiende algunos resultados conocidos. Finalmente introducimos una nueva clase de polinomios tipo Apostol generalizados basados en los polinomios de Hermite de dos variables ${_{H}}\mathcal{Q}_n^{[m-1,\alpha]}(x,b,c;\lambda;u,v)$ y mencionamos propiedades b�sicas.


\vspace{2cm}

\noindent \textbf{Palabras Claves:} polinomios tipo Apostol generalizados; relaci�n de recurrencia; ecuaci�n diferencial; polinomios de Apostol-Bernoulli; polinomios de Genocchi; polinomios de Jacobi; polinomios de Hermite; polinomios de Laguerre; polinomios de Charlier; polinomios de Bessel; polinomios de Bernoulli generalizados; n�meros de Stirling de segunda clase.




\renewcommand\indexname{\'INDICE ALFAB\'ETICO}
\renewcommand\bibname{REFERENCIAS BIBLIOGR\'AFICAS}
\renewcommand\contentsname{\'INDICE GENERAL}
\tableofcontents
\markboth{\'INDICE GENERAL}{\'INDICE GENERAL}
\chapter*{LISTA DE S\'IMBOLOS}

\addcontentsline{toc}{chapter}{LISTA DE S\'IMBOLOS}


\vspace{1.0cm}

\begin{tabular}{lll}
$\ZZ, \NN, \RR, \CC $ & & n\'umeros enteros, naturales, reales y complejos respectivamente.\\
$\mathcal{P}$ & & espacio vectorial de los polinomios con coeficientes complejos.\\
$\mathcal{P}_n$& & subespacio vectorial de los polinomios de grado a lo sumo $n$.\\
$\langle\,\cdot,\cdot\,\rangle_{\mu}$ & & producto interno asociado a la medida $\mu$.\\
$\mu$& &medida de Borel.\\
$(X,\beta,\mu)$& & espacio de medida.\\
$\sop{\mu}\;$ & & sorporte de la medida $\mu$.\\
$H_n(x)$ & & $n$-\'esimo polinomio de Hermite.\\
$L_n^{(\alpha)}(x)$ & & $n$-\'esimo polinomio de Laguerre.\\
$P_n^{(\alpha,\beta)}(x)$ & & $n$-\'esimo polinomio de Jacobi.\\
$\mathfrak{B}_n(\lambda;x)$ & & $n$-\'esimo polinomio de Apostol-Bernoulli.\\
$\mathfrak{E}_n(\lambda;x)$ & & $n$-\'esimo polinomio de Apostol-Euler.\\
$\mathfrak{G}_n(\lambda;x)$ & & $n$-\'esimo polinomio de Apostol-Genocchi.\\
$S(m,n)$ & & n\'umeros de Stirling de primera clase.\\
$\mathfrak{S}(n,k)$ & & n\'umeros de Stirling de segunda clase.\\
$C_n^{(a)}(x)$ & & $n$-\'esimo polinomio de Charlier.\\
$\varphi_n (x,c)$ & & $n$-\'esimo polinomio de Bessel generalizado.\\
$\mathcal{F}_n^{[m-1,\alpha]}(x;\lambda;u,v)$ & & $n$-\'esimo polinomio tipo Apostol generalizado.\\



\end{tabular}














\pagenumbering{arabic}
\setcounter{page}{1}
\mainmatter

\chapter*{\textcolor{blue}{Introducci�n}}

\addstarredchapter{Introducci�n}

%\iflettrine


\lettrine[lines=4,loversize=-0.1,lraise=0.1,lhang=.2]{A}{}{l interior de los experimentos estad\'{i}sticos la teor\'{i}a de
los dise�os �ptimos ha sido desarrollada. En general el tema de esta
teor\'{i}a es que para un apropiado modelo, si queremos poner �nfasis
sobre una cualidad particular de los par�metros a estimar, entonces
la configuraci�n experimental deber\'{i}a ser elegida de acuerdo a
ciertos criterios con sentido estad\'{i}stico. En la literatura relacionada
con los dise�os �ptimos, un prominente autor fue Kiefer (1959), el
cu�l present� los principales conceptos, tales como dise�os aproximados
y una variedad de criterios de �ptimalidad para esta rama de los dise�os
de experimentos; Kiefer, en particular dio el nombre $D-$optimalidad
al criterio introducido por Wald (1943), este criterio es el m�s comunmente
aplicado y est� definido en funci�n del determinante de la matriz
de covarianza.\\
}

M�s recientemente son reconocidos los libros de Atkinson y Donev (1992)
y Pukelsheim (1993), donde los autores hacen una presentaci�n estad\'{i}stica
formal de los dise�os �ptimos. El presente trabajo se ha organizado
en tres cap\'{i}tulos: el cap\'{i}tulo uno (Preliminares) contiene
conceptos generales que sirven de apoyo y base a la teor\'{i}a que
se desarrolla en los siguientes dos cap\'{i}tulos. El cap\'{i}tulo
dos trata sobre dise�os �ptimos en la presencia de efectos de bloques
aleatorios, y en el cap\'{i}tulo tres se daran a conocer las conclusiones
y una serie de problemas abiertos para futuras investigaciones relacionadas
con el tema central de este trabajo de investigaci�n. 

\include{apostolB1}
 \chapter{SOBRE LAS DISTINTAS GENERALIZACIONES DE POLINOMIOS  DE APOSTOL-BERNOULLI, DE APOSTOL-EULER Y DE APOSTOL-GENOCCHI}\markboth{CAP\'ITULO 2\, GENERALIZACIONES DE LOS POL. DE A-B, A-E Y A-G}{CAP\'ITULO 2\, GENERALIZACIONES DE LOS POL. DE A-B, DE A-E Y A-G}




%recuperar numeracion arabica%
\def\thechapter{\arabic{chapter}}

Consideremos los polinomios de Apostol-Bernoulli $\mathfrak{B}_n(x;\lambda)$, los polinomios de Apostol-Euler $\mathfrak{E}_n(x;\lambda)$ y los polinomios de Apostol-Genocchi $\mathfrak{G}_n(x;\lambda)$, definidos en (\ref{eq3}), (\ref{eq300}) y (\ref{eq301}) respectivamente. En este cap\'{\i}tulo mostraremos, las diferentes generalizaciones que han tenido estos polinomios hasta llegar a una de las m�s actuales conocidas. Lo anterior implica que estudiaremos fundamentalmente las definiciones y algunas propiedades de los polinomios de Apostol-Euler generalizados, los polinomios de Apostol-Genocchi generalizados, las nuevas familias y clases de polinomios de Apostol-Bernoulli, de Apostol-Euler y de Apostol-Genocchi generalizados y finalmente los polinomios tipo Apostol generalizados. 

\section{Polinomios de Apostol-Bernoulli, de Apostol-Euler y de Apostol-Genocchi generalizados}
Muchas interesantes generalizaciones de los polinomios de Apostol-Bernoulli, Apostol-Euler y Apostol-Genocchi han sido dadas. En part�cular Luo y Srivastava \cite{gen, REF9} introdujeron los polinomios de Apostol-Bernoulli generalizados
$\mathfrak{B}_{n}^{(\alpha)}(x;\lambda)$ de orden $\alpha \in \CC$; Luo \cite{REF6} invent� los polinomios de Apostol-Euler $\mathfrak{E}_{n}^{(\alpha)}(x;\lambda)$ de orden $\alpha \in \CC$ y los polinomios de Genocchi $\mathfrak{G}_{n}^{(\alpha)}(x;\lambda)$ de orden $\alpha \in \CC$ en \cite{Genpolt}

\begin{definition}\label{abp}
Los polinomios de Apostol-Bernoulli generalizados $\mathfrak{B}_{n}^{(\alpha)}(x;\lambda)$, los polinomios de Apostol Euler generalizados $\mathfrak{E}_{n}^{(\alpha)}(x;\lambda)$ y los polinomios de Apostol-Genocchi generalizados $\mathfrak{G}_{n}^{(\alpha)}(x;\lambda)$ de orden $\alpha \in \CC$ est�n definidos por las siguientes funciones generatrices:
\begin{equation}\label{eqae}
\displaystyle\left(\frac{z}{\lambda e^z-1}\right)^{\alpha}e^{xz} =\displaystyle \sum\limits_{n=0}^{\infty}
\mathfrak{B}_{n}^{(\alpha)}(x;\lambda)\frac{z^n}{n!}; \quad (|z|<|\log(\lambda)|),
\end{equation}

\begin{equation}\label{eqae1}
\displaystyle\left(\frac{2}{\lambda e^z+1}\right)^{\alpha}e^{xz} =\displaystyle \sum\limits_{n=0}^{\infty}
\mathfrak{E}_{n}^{(\alpha)}(x;\lambda)\frac{z^n}{n!}; \quad (|t|<|\log(-\lambda)|),
\end{equation}

\begin{equation}\label{eqae2}
\displaystyle\left(\frac{2z}{\lambda e^z+1}\right)^{\alpha}e^{xz} =\displaystyle \sum\limits_{n=0}^{\infty}
\mathfrak{G}_{n}^{(\alpha)}(x;\lambda)\frac{z^n}{n!}; \quad (|z|<|\log(-\lambda)|).
\end{equation}
\end{definition}

\vspace{0.25cm}

\begin{obs}
Los polinomios de Apostol-Bernoulli $\mathfrak{B}_{n}(x;\lambda) = \mathfrak{B}_{n}^{(1)}(x;\lambda)$ y los polinomios de Apostol-Bernoulli generalizados $B_{n}^{(\alpha)}(x) = \mathfrak{B}_{n}^{(\alpha)}(x;1)$.\\
Los polinomios de Apostol-Euler $\mathfrak{E}_{n}(x;\lambda) = \mathfrak{E}_{n}^{(1)}(x;\lambda)$ y los polinomios de Apostol-Euler generalizados $E_{n}^{(\alpha)}(x) = \mathfrak{E}_{n}^{(\alpha)}(x;1)$.\\
Los polinomios de Apostol-Genocchi $\mathfrak{G}_{n}(x;\lambda) = \mathfrak{G}_{n}^{(1)}(x;\lambda)$ y los polinomios de Apostol-Genocchi generalizados $G_{n}^{(\alpha)}(x) = \mathfrak{G}_{n}^{(\alpha)}(x;1)$.
\end{obs}


\section{Nuevas familias de polinomios de Apostol-Euler y de Apostol-Genocchi generalizados}
En este apartado presentaremos dos familias de polinomios de Apostol-Euler y Apostol-Genocchi generalizados  que fueron introducidas por Srivastava et. al.\cite{Srivasnew}. Ellos investigaron las siguientes formas.

\vspace{0.25cm}
\begin{definition}\label{abs}
Sean $a,b,c \in \RR^{+} \quad (a\neq b)$ y $n \in \NN_0:=\NN \cup \{0\}$. Entonces los polinomios de Apostol-Euler generalizados $\mathfrak{E}_n^{(\alpha)}(x;\lambda;a,b,c)$ de orden $\alpha \in \CC$ est�n definidos por la siguiente funci�n generatriz:
\begin{equation}\label{eqaps}
\displaystyle\left(\frac{2}{\lambda b^z+a^z}\right)^{\alpha}c^{xz} =\sum\limits_{n=0}^{\infty}
\mathfrak{E}_n^{(\alpha)}(x;\lambda;a,b,c)\frac{z^n}{n!} 
\end{equation}

\begin{center}
$\displaystyle\quad \left(\left|z\log \left(\frac{b}{a}\right)\right|<|\log(-\lambda)|; 1^\alpha:=1; x\in \RR\right)$
\end{center}
\end{definition}

\vspace{0.25cm}

\begin{definition}\label{abs1}
Sean $a,b,c \in \RR^{+} \quad (a\neq b)$ y $n \in \NN_0:=\NN \cup \{0\}$. Entonces los polinomios Apostol-Genocchi generalizados $\mathfrak{G}_n^{(\alpha)}(x;\lambda;a,b,c)$ de orden $\alpha \in \CC$ est�n definidos por la siguiente funci�n generatriz:
\begin{equation}\label{eqaps1}
\displaystyle\left(\frac{2z}{\lambda b^z+a^z}\right)^{\alpha}c^{xz} =\sum\limits_{n=0}^{\infty}
\mathfrak{G}_n^{(\alpha)}(x;\lambda;a,b,c)\frac{z^n}{n!} 
\end{equation}

\begin{center}
$\displaystyle\quad \left(\left|z\log \left(\frac{b}{a}\right)\right|<|\log(-\lambda)|; 1^\alpha:=1; x\in \RR\right)$
\end{center}
\end{definition}

\vspace{0.25cm}

\section{Nuevas clases de polinomios de Apostol-Bernoulli, de Apostol-Euler y de Apostol-Genocchi generalizados}
En esta secci�n ilustraremos algunas nuevas clases de polinomios de Apostol-Bernoulli generalizados $\mathcal{B}_n^{[m-1,\alpha]}(x,b,c;\lambda)$ que fueron introducidos por Tremblay, Gaboury y Fug\`{e}re \cite{Tremblay1} y nuevas clases de polinomios de Apostol-Euler $\mathcal{E}_n^{[m-1,\alpha]}(x,b,c;\lambda)$ y Apostol-Genocchi generalizados $\mathcal{G}_n^{[m-1,\alpha]}(x,b,c;\lambda)$ que fueron introducidos por los mismos autores en \cite{Tremblay2}.\\

\begin{definition}\label{Gabory}
Para par�metros complejos o reales arbitrarios $\alpha$ y para $b,c \in \RR^{+}$, la nueva clase de polinomios Apostol-Bernoulli generalizados $\mathcal{B}_n^{[m-1,\alpha]}(x,b,c;\lambda)$, la nueva clase de polinomios Apostol-Euler generalizados $\mathcal{E}_n^{[m-1,\alpha]}(x,b,c;\lambda)$ y la nueva clase de polinomios Apostol-Genocchi generalizados $\mathcal{G}_n^{[m-1,\alpha]}(x,b,c;\lambda)$, $m \in \NN$, $\lambda \in \CC$ est�n definidos en un entorno adecuado de $z=0$, a trav�s de la siguientes funciones generatrices:

\begin{equation}\label{eq5}
\displaystyle\left(\frac{z^{m}}{\lambda b^z-\sum\limits_{l=0}^{m-1}\frac{(z\log b)^l}{l!}}\right)^{\alpha}c^{xz} =\displaystyle\sum\limits_{n=0}^{\infty}
\mathcal{B}_n^{[m-1,\alpha]}(x,b,c;\lambda)\frac{z^n}{n!},
\end{equation}

\begin{equation}\label{eq51}
\displaystyle\left(\frac{2^{m}}{\lambda b^z+\sum\limits_{l=0}^{m-1}\frac{(z\log b)^l}{l!}}\right)^{\alpha}c^{xz} =\displaystyle\sum\limits_{n=0}^{\infty}
\mathcal{E}_n^{[m-1,\alpha]}(x,b,c;\lambda)\frac{z^n}{n!},
\end{equation}

\begin{equation}\label{eq52}
\displaystyle\left(\frac{(2z)^{m}}{\lambda b^z+\sum\limits_{l=0}^{m-1}\frac{(z\log b)^l}{l!}}\right)^{\alpha}c^{xz} =\displaystyle\sum\limits_{n=0}^{\infty}
\mathcal{G}_n^{[m-1,\alpha]}(x,b,c;\lambda)\frac{zs^n}{n!}.
\end{equation}

\end{definition}

\vspace{0.25cm}

\begin{obs}
Si hacemos $b=c=e$ en (\ref{eq5}), (\ref{eq51}) y (\ref{eq52}) obtenemos

\begin{equation}\label{eq50}
\displaystyle\left(\frac{z^{m}}{\lambda e^z-\sum\limits_{l=0}^{m-1}\frac{z^l}{l!}}\right)^{\alpha}e^{xz} =\displaystyle\sum\limits_{n=0}^{\infty}
\mathcal{B}_n^{[m-1,\alpha]}(x;\lambda)\frac{z^n}{n!},
\end{equation}

\begin{equation}\label{eq511}
\displaystyle\left(\frac{2^{m}}{\lambda e^z+\sum\limits_{l=0}^{m-1}\frac{z^l}{l!}}\right)^{\alpha}e^{xz} =\displaystyle\sum\limits_{n=0}^{\infty}
\mathcal{E}_n^{[m-1,\alpha]}(x;\lambda)\frac{z^n}{n!},
\end{equation}

\begin{equation}\label{eq522}
\displaystyle\left(\frac{(2z)^{m}}{\lambda e^z+\sum\limits_{l=0}^{m-1}\frac{z^l}{l!}}\right)^{\alpha}e^{xz} =\displaystyle\sum\limits_{n=0}^{\infty}
\mathcal{G}_n^{[m-1,\alpha]}(x;\lambda)\frac{z^n}{n!}.
\end{equation}

\end{obs}

\vspace{0.25cm}

\begin{obs}
Si hacemos $m=1$ en (\ref{eq50}), (\ref{eq511}) y (\ref{eq522}) obtenemos (\ref{eqae}), (\ref{eqae1}) y (\ref{eqae2}) respectivamente.
\end{obs}

\vspace{0.25cm}

\begin{lemma}
 Los polinomios de Bernoulli generalizados $B_{n}^{[m-1]}(x):=\mathcal{B}_{n}^{[m-1,1]}(x;1)$, $m \in \NN$ verifican la siguiente relaci�n:
\begin{equation}\label{eq60}
x^n=\displaystyle\sum\limits_{k=0}^{n}\binom{n}{k}\frac{k!}{(k+m)!}B_{n-k}^{[m-1]}(x).
\end{equation}
\end{lemma}

\vspace{0.25cm}

Para detalles de la prueba del anterior resultados, ver \cite{REF_14}, p.158, Ec.(2.6).


\vspace{0.25cm}

\section{Polinomios tipo Apostol generalizados}
Consideramos en esta secci�n la generalizaci�n m�s importante y m�s actual para los polinomios de Apostol-Bernoulli, Apostol-Euler y Apostol-Gencocchi, la cual fue introducida en 2013 por Lu y Luo en \cite{atp1}.\\

\begin{definition}\label{abp1}
Los polinomios tipo Apostol generalizados $\mathcal{F}_n^{(\alpha)}(x;\lambda;u,v) \quad \alpha \in \NN_0,\lambda, u, v \in \CC)$ de orden $\alpha$ est�n definidos a trav�s de la siguiente funci�n generatriz:
\begin{equation}\label{eqapt1}
\displaystyle\left(\frac{2^{u}z^{v}}{\lambda e^z+1}\right)^{\alpha}e^{xz} =\sum\limits_{n=0}^{\infty}
\mathcal{F}_n^{(\alpha)}(x;\lambda;u,v)\frac{z^n}{n!}; \quad (|z|<|log(-\lambda)|),
\end{equation}
donde
\begin{equation}
\mathcal{F}_n^{(\alpha)}(\lambda;u,v)=:\mathcal{F}_n^{(\alpha)}(0;\lambda;u,v) 
\end{equation}
denotan los tambi�n llamados n�meros de tipo Apostol de orden $\alpha$.
\end{definition}

\vspace{0.25cm}

So that, by comparing \textit{Definition \ref{abp1}} and \textit{Definition \ref{abp}} 
\begin{align}
& \mathfrak{B}_n^{(\alpha)}(x;\lambda)=(-1)^{\alpha}\mathcal{F}_n^{(\alpha)}(x;-\lambda;0,1),\\
& \mathfrak{E}_n^{(\alpha)}(x;\lambda)=\mathcal{F}_n^{(\alpha)}(x;\lambda;1,0),\\
& \mathfrak{G}_n^{(\alpha)}(x;\lambda)=\mathcal{F}_n^{(\alpha)}(x;\lambda;1,1).
\end{align}

\vspace{0.25cm}


 
\chapter{UNA NUEVA CLASE DE POLINOMIOS TIPO APOSTOL GENERALIZADOS}
\markboth{CAP\'ITULO 3\, UNA NUEVA CLASE DE POL. TIPO APOSTOL GENERALIZADOS}{CAP\'ITULO 3\quad UNA NUEVA CLASE DE POL. TIPO APOSTOL GENERALIZADOS}
\label{AB3}

%recuperar numeracion arabica%
\def\thechapter{\arabic{chapter}}
En este cap�tulo se introducir\'a una nueva clase de polinomios tipo Apostol generalizados denotada por $\mathcal{Q}_n^{[m-1,\alpha]}(x,b,c;\lambda;u,v)$, as\'{\i} como tambi\'en, se estudiar\'an sus propiedades b�sicas incluidas su f�rmula de recurrencia y ecuaci�n diferencial que satisfacen. Para finalizar, se establecer�n f�rmulas de conexi�n entre $\mathcal{Q}_n^{[m-1,\alpha]}(x,b,c;\lambda;u,v)$ y los polinomios ortogonales cl�sicos, con los polinomios de Genocchi, Charlier, Bessel, Bernoulli generalizados y los n�meros de Stirling de segunda clase.\\


La siguiente definici�n que generaliza la \textit{Definici�n \ref{Gabory}} en una sola expresi�n y adem�s proporciona una extensi�n natural para los polinomios tipo Apostol generalizados $\mathcal{F}_n^{(\alpha)}(x;\lambda;u,v) \quad (\alpha \in \NN_0,\lambda, u, v \in \CC)$ de orden $\alpha$ definidos en (\ref{eqapt1}).


\medskip

\begin{definition}\label{abp14}
La nueva clase para los polomios tipo Apostol generalizados $\mathcal{Q}_n^{[m-1,\alpha]}(x,b,c;\lambda;u,v)$ con $(\alpha, u, v\in \CC$ y $b, c \in \RR^{+})$ de orden $\alpha$ est�n definidos por la siguiente funci�n generatriz:
\begin{equation}\label{newpol}
\displaystyle\left(\frac{(2^ut^v)^{m}}{\lambda b^t+\sum\limits_{l=0}^{m-1}\frac{(t\log b)^l}{l!}}\right)^{\alpha}c^{xt} =\displaystyle\sum\limits_{n=0}^{\infty}
\mathcal{Q}_n^{[m-1,\alpha]}(x,b,c;\lambda;u,v)\frac{t^n}{n!}; \quad |t\log b|<|log(-\lambda)|,
\end{equation}
donde
\begin{equation}\label{newnumbers}
\mathcal{Q}_n^{[m-1,\alpha]}(b,c;\lambda;u,v):=\mathcal{Q}_n^{[m-1,\alpha]}(0,b,c;\lambda;u,v)
\end{equation}
denotan los tambi�n llamados la nueva clase de n�meros de tipo Apostol generalizados de orden $\alpha$.
\end{definition}

\vspace{0.25cm}

As� que, comparando la \textit{Definici�n \ref{abp14}} y la \textit{Definici�n \ref{Gabory}} tenemos: 
\begin{align}
& \mathcal{B}_n^{[m-1,\alpha]}(x;b,c;\lambda)=(-1)^{\alpha}\mathcal{Q}_n^{[m-1,\alpha]}(x;b,c;-\lambda;0,1),\\
& \mathcal{E}_n^{[m-1,\alpha]}(x;b,c;\lambda)=\mathcal{Q}_n^{[m-1,\alpha]}(x;b,c;\lambda;1,0),\\
& \mathcal{G}_n^{[m-1,\alpha]}(x;b,c;\lambda)=\mathcal{Q}_n^{[m-1,\alpha]}(x;b,c;\lambda;1,1).
\end{align}

\vspace{0.25cm}

\begin{obs}
$\mathcal{Q}_n^{(\alpha)}(x;\lambda;u,v):=\mathcal{Q}_n^{[0,\alpha]}(x,e,e;\lambda;u,v)$
\end{obs}

\vspace{0.25cm}

Analogamente, usando las ideas de la \textit{Definici�n \ref{abp1}}, podemos escribir las \textit{Definiciones \ref{abs} and \ref{abs1}} en una sola expresi�n, como se sigue:

\vspace{0.25cm}

\begin{definition}\label{newpol2}
Sean $a,b,c \in \RR^{+} \quad (a\neq b)$, $u, v \in \CC$ y $n \in \NN_0:=\NN \cup \{0\}$. Entonces la nueva clase de polinomios tipo Apostol basado en Srivastava $\mathcal{Q}_n^{(\alpha)}(x;\lambda;a,b,c;u,v)$ de orden $\alpha \in \CC$ est� definida a trav�s de la siguiente funci�n generatriz:
\begin{equation}\label{eqnewpol2}
\displaystyle\left(\frac{2^uz^v}{\lambda b^z+a^z}\right)^{\alpha}c^{xz} =\sum\limits_{n=0}^{\infty}
\mathcal{Q}_n^{(\alpha)}(x;\lambda;a,b,c;u,v)\frac{z^n}{n!} 
\end{equation}

\begin{center}
$\displaystyle\quad \left(\left|z\log \left(\frac{b}{a}\right)\right|<|\log(-\lambda)|; 1^\alpha:=1; x\in \RR\right)$
\end{center}
\end{definition}

\vspace{0.25cm}

Los nuevos polinomios en (\ref{eqnewpol2}) tambi�n son una nueva clase de polinomios tipo Apostol al igual los definidos en (\ref{newpol}). Estos verifican lo siguiente:

\begin{align}
& \mathfrak{B}_n^{(\alpha)}(x;\lambda;a,b,c):=(-1)^{\alpha}\mathcal{Q}_n^{(\alpha)}(x;-\lambda;a,b,c;0,1),\\
& \mathfrak{E}_n^{(\alpha)}(x;\lambda;a,b,c)=\mathcal{Q}_n^{(\alpha)}(x;\lambda;a,b,c;1,0),\\
& \mathfrak{G}_n^{(\alpha)}(x;\lambda;a,b,c)=\mathcal{Q}_n^{(\alpha)}(x;\lambda;a,b,c;1,1).
\end{align}


\vspace{0.5cm}
\section{Algunas propiedades b�sicas de $\mathcal{Q}_n^{[m-1,\alpha]}(x,b,c;\lambda;u,v)$}\label{3d1}
En est� secci�n estudiaremos algunas de las propiedades b�sicas de los polinomios $\mathcal{Q}_n^{[m-1,\alpha]}(x,b,c;\lambda;u,v)$.

\medskip

\begin{teo}\label{3cau1}
La nueva clase de polinomios tipo Apostol generalizados $\mathcal{Q}_n^{[m-1,\alpha]}(x,b,c;\lambda;u,v)$ satisface las siguientes relaciones:
\begin{equation}\label{eq8}
\mathcal{Q}_n^{[m-1,\alpha+\beta]}(x+y,b,c;\lambda;u,v)=\displaystyle\sum\limits_{j=0}^{n}
\binom{n}{j}\mathcal{Q}_j^{[m-1,\alpha]}(x,b,c;\lambda;u,v)\mathcal{Q}_{n-j}^{[m-1,\beta]}(y,b,c;\lambda;u,v),
\end{equation}
\begin{equation}\label{eq9}
\mathcal{Q}_n^{[m-1,\alpha]}(x+y,b,c;\lambda;u,v)=\displaystyle\sum\limits_{j=0}^{n}
\binom{n}{j}\mathcal{Q}_j^{[m-1,\alpha]}(y,b,c;\lambda;u,v)(x\log c)^{n-j}.
\end{equation}
\end{teo}

\begin{proof}
 Por (\ref{newpol}), tenemos:

\begin{eqnarray*}
\displaystyle \sum\limits_{n=0}^{\infty} \mathcal{Q}_n^{[m-1,\alpha+\beta]}(x+y,b,c;\lambda;u,v)\frac{t^n}{n!} &=& \displaystyle\left(\frac{(2^ut^v)^{m}}{\lambda b^t+\sum\limits_{l=0}^{m-1}\frac{(t\log b)^l}{l!}}\right)^{\alpha+\beta}c^{(x+y)t}\\
&=& \displaystyle\left(\frac{(2^ut^v)^{m}}{\lambda b^t+\sum\limits_{l=0}^{m-1}\frac{(t\log b)^l}{l!}}\right)^{\alpha}c^{xt}\left(\frac{(2^ut^v)^{m}}{\lambda b^t+\sum\limits_{l=0}^{m-1}\frac{(t\log b)^l}{l!}}\right)^{\beta}c^{yt}\\
&=& \displaystyle\sum\limits_{n=0}^{\infty}
\mathcal{Q}_n^{[m-1,\alpha]}(x,b,c;\lambda;u,v)\frac{t^n}{n!}\sum\limits_{n=0}^{\infty}
\mathcal{Q}_n^{[m-1,\beta]}(y,b,c;\lambda;u,v)\frac{t^n}{n!}\\
&=& \displaystyle \sum\limits_{n=0}^{\infty}\sum\limits_{k=0}^{n}\binom{n}{k}\mathcal{Q}_k^{[m-1,\alpha]}(x,b,c;\lambda;u,v)\mathcal{Q}_{n-k}^{[m-1,\beta]}(y,b,c;\lambda;u,v)\frac{t^n}{n!}.\\
\end{eqnarray*}
Comparando los coeficientes de $\frac{t^n}{n!}$ en ambos miembros de la ecuaci�n anterior, llegamos a (\ref{eq8}) rapidamente.

\vspace{0.25cm}

Haciendo $\beta = 0$ en (\ref{eq8}) tenemos:
\begin{eqnarray*}
\displaystyle \sum\limits_{n=0}^{\infty} \mathcal{Q}_n^{[m-1,\alpha]}(x+y,b,c;\lambda;u,v)\frac{t^n}{n!} &=& \displaystyle\left(\frac{(2^ut^v)^{m}}{\lambda b^t+\sum\limits_{l=0}^{m-1}\frac{(t\log b)^l}{l!}}\right)^{\alpha}c^{(x+y)t}\\
&=& \displaystyle\left(\frac{(2^ut^v)^{m}}{\lambda b^t+\sum\limits_{l=0}^{m-1}\frac{(t\log b)^l}{l!}}\right)^{\alpha}c^{yt}c^{xt}\\
&=& \displaystyle\left(\frac{(2^ut^v)^{m}}{\lambda b^t+\sum\limits_{l=0}^{m-1}\frac{(t\log b)^l}{l!}}\right)^{\alpha}c^{yt}e^{\log c^{xt}}\\
&=& \displaystyle\left(\frac{(2^ut^v)^{m}}{\lambda b^t+\sum\limits_{l=0}^{m-1}\frac{(t\log b)^l}{l!}}\right)^{\alpha}c^{yt}e^{(x\log c)t}\\
&=& \displaystyle\sum\limits_{n=0}^{\infty}
\mathcal{Q}_n^{[m-1,\alpha]}(y,b,c;\lambda;u,v)\frac{t^n}{n!}\sum\limits_{n=0}^{\infty}
(x\log c)^n\frac{t^n}{n!}\\
&=& \displaystyle \sum\limits_{n=0}^{\infty}\sum\limits_{k=0}^{n}\binom{n}{k}\mathcal{Q}_k^{[m-1,\alpha]}(y,b,c;\lambda;u,v)(x\log c)^{n-k}\frac{t^n}{n!}.
\end{eqnarray*}

Comparando los coeficientes de $\frac{t^n}{n!}$ en ambos miembros de la ecuaci�n anterior, llegamos a (\ref{eq9}).
\end{proof}

\vspace{0.25cm}

\begin{coro}\label{3cau12}
La nueva clase de polinomios tipo Apostol generalizados $\mathcal{Q}_n^{[m-1,\alpha]}(x,b,c;\lambda;u,v)$ satisfacen las siguientes relaciones:
\begin{align}
\label{eq823}
\mathcal{Q}_n^{[m-1,\alpha]}(x,b,c;\lambda;u,v)= & \displaystyle\sum\limits_{j=0}^{n}
\binom{n}{j}\mathcal{Q}_{n-j}^{[m-1,\alpha-1]}(b,c;\lambda;u,v)\mathcal{Q}_j^{[m-1]}(x,b,c;\lambda;u,v),\\
\label{eq349}
\mathcal{Q}_n^{[m-1,\alpha]}(x,b,c;\lambda;u,v)= & \displaystyle\sum\limits_{j=0}^{n}
\binom{n}{j}\mathcal{Q}_{n-j}^{[m-1,\alpha]}(b,c;\lambda;u,v)(x\log c)^{j},
\end{align}
donde $\mathcal{Q}_{n}^{[m-1,\alpha]}(b,c;\lambda;u,v)$ son los tambi�n llamados nueva clase de n�meros tipo Apostol generalizados de orden $\alpha$ definidos por (\ref{newnumbers}).
\end{coro}

\medskip

\begin{proof}
Haciendo $\alpha=1$, $\beta=\alpha-1$ y $y=0$ en la Ec.(\ref{eq8}) del \textit{Teorema \ref{3cau1}}, tenemos (\ref{eq823}). Cuando $y=0$ en la Ec.(\ref{eq9}) of the \textit{Teorema \ref{3cau1}} y aplicando (\ref{s4}) podemos obtener (\ref{eq349}).
\end{proof}

\medskip
\begin{obs}
Tomando $m=1$ y $b=c=e$ en el \textit{Corolario \ref{3cau12}}, obtenemos:
\begin{equation}\label{eq8231}
\mathcal{F}_n^{(\alpha)}(x;\lambda;u,v)=\displaystyle\sum\limits_{j=0}^{n}
\binom{n}{j}\mathcal{F}_{n-j}^{(\alpha-1)}(\lambda;u,v)\mathcal{F}_j(x;\lambda;u,v),
\end{equation}
la cual corresponde a la Ec.(3.5) de \cite{atp1},
\begin{equation}\label{eq3491}
\mathcal{F}_n^{(\alpha)}(x;\lambda;u,v)=\displaystyle\sum\limits_{j=0}^{n}
\binom{n}{j}\mathcal{F}_{n-j}^{(\alpha)}(\lambda;u,v)x^{j},
\end{equation}
la cual corresponde a la Ec.(3.1) de \cite{atp1}.
\end{obs}

\medskip

\begin{teo}\label{par}
La nueva clase de polinomios tipo Apostol generalizados $\mathcal{Q}_n^{[m-1,\alpha]}(x,b,c;\lambda;u,v)$ satisfacen la siguiente relaci�n:
\begin{equation}\label{e7}
\displaystyle \frac{\partial^p}{\partial x^p} \mathcal{Q}_n^{[m-1,\alpha]}(x,b,c;\lambda;u,v)= \frac{n!}{(n-p)!}(\log c)^{p}\mathcal{Q}_{n-p}^{[m-1,\alpha]}(x,b,c;\lambda;u,v).
\end{equation}
\end{teo}

\begin{proof}
Usando la ecuaci�n (\ref{eq349}), tenemos\\
\begin{equation}\nonumber
\mathcal{Q}_n^{[m-1,\alpha]}(x,b,c;\lambda;u,v)=\displaystyle\sum\limits_{k=0}^{n}
\binom{n}{k}\mathcal{Q}_k^{[m-1,\alpha]}(b,c;\lambda;u,v)(x\log c)^{n-k}
\end{equation}
diferenciando parcialmente $p$ veces respecto a la variable $x$ en ambos miembros, se sigue que:
\begin{eqnarray*}
\displaystyle \frac{\partial^p}{\partial x^p}\mathcal{Q}_n^{[m-1,\alpha]}(x,b,c;\lambda;u,v) &=& \frac{\partial^p}{\partial x^p}\left[\displaystyle \sum\limits_{k=0}^{n}\binom{n}{k}\mathcal{Q}_k^{[m-1,\alpha]}(b,c;\lambda;u,v)\displaystyle (x\log c)^{n-k}\right]\\
&=& \frac{\partial^p}{\partial x^p}\left[\displaystyle \sum\limits_{k=0}^{n}\binom{n}{k}\mathcal{Q}_k^{[m-1,\alpha]}(b,c;\lambda;u,v)\displaystyle (\log c)^{n-k}x^{n-k}\right]\\
&=& \displaystyle \sum\limits_{k=0}^{n-p}\binom{n}{k}\mathcal{Q}_k^{[m-1,\alpha]}(b,c;\lambda;u,v)(\log c)^{n-k} \frac{\partial^p}{\partial x^p}\left(x^{n-k}\right)\\
&=& \displaystyle \sum\limits_{k=0}^{n-p}\binom{n}{k}\mathcal{Q}_k^{[m-1,\alpha]}(b,c;\lambda;u,v)(\log c)^{n-k} \frac{(n-k)!}{(n-p-k)!} x^{n-p-k}\\
\end{eqnarray*}
desarrollando el combinatorio, obtenemos
\begin{eqnarray*}
\displaystyle \frac{\partial^p}{\partial x^p}\mathcal{Q}_n^{[m-1,\alpha]}(x,b,c;\lambda;u,v) &=& \displaystyle \sum\limits_{k=0}^{n-p}\mathcal{Q}_k^{[m-1,\alpha]}(b,c;\lambda;u,v)(\log c)^{n-k} \frac{(n-k)!}{(n-p-k)!}\frac{n!}{k!(n-k)!} x^{n-p-k}\\
&=& \displaystyle \sum\limits_{k=0}^{n-p}\mathcal{Q}_k^{[m-1,\alpha]}(b,c;\lambda;u,v)(\log c)^{n-k} \frac{(n-p)!}{k!(n-p-k)!}\frac{n!}{(n-p)!} x^{n-p-k}\\
&=& \displaystyle \frac{n!}{(n-p)!}\sum\limits_{k=0}^{n-p}\binom{n-p}{k}\mathcal{Q}_k^{[m-1,\alpha]}(b,c;\lambda;u,v)(\log c)^{p} (x\log c)^{n-p-k}\\
&=& \displaystyle \frac{n!}{(n-p)!}(\log c)^{p}\sum\limits_{k=0}^{n-p}\binom{n-p}{k}\mathcal{Q}_k^{[m-1,\alpha]}(b,c;\lambda;u,v) (x\log c)^{n-p-k}\\
&=& \displaystyle \frac{n!}{(n-p)!}(\log c)^{p}\mathcal{Q}_{n-p}^{[m-1,\alpha]}(x,b,c;\lambda;u,v).
\end{eqnarray*}
\end{proof}

\vspace{0.25cm}

Al establecer $m=1$ y $b=c=e$ en el \textit{Teorema \ref{par}}, tenemos el siguiente corolario.

\vspace{0.25cm}

\begin{coro}\label{par23}
\begin{equation}\label{e723}
\displaystyle \frac{\partial^p}{\partial x^p} \mathcal{F}_n^{(\alpha)}(x;\lambda;u,v)= \frac{n!}{(n-p)!}\mathcal{F}_{n-p}^{(\alpha)}(x;\lambda;u,v).
\end{equation}
la cual es justamente la Ec.(2.10) de \cite{atp1}. 
\end{coro}

\vspace{0.25cm}

\begin{teo}\label{par2}
La nueva clase de polinomios tipo Apostol generalizados $\mathcal{Q}_n^{[m-1,\alpha]}(x,b,c;\lambda;u,v)$ satisfacen la siguiente relaci�n:
\begin{equation}\label{e72}
t^p\left(\frac{(2^ut^v)^{m}}{\lambda b^t+\sum\limits_{l=0}^{m-1}\frac{(t\log b)^l}{l!}}\right)^{\alpha}c^{xt} =\sum\limits_{n=0}^{\infty}\frac{n!}{(n-p)!}\mathcal{Q}_{n-p}^{[m-1,\alpha]}(x,b,c;\lambda;u,v)\frac{t^n}{n!}.
\end{equation}
\end{teo}

\vspace{0.25cm}

\begin{proof}
Consideremos (\ref{newpol}) y derivemos parcialmente $p$ veces en ambos miembros respecto a la variable $x$, entonces tenemos:
\begin{align*}
t^p(\log c)^p\left(\frac{(2^ut^v)^{m}}{\lambda b^t+\sum\limits_{l=0}^{m-1}\frac{(t\log b)^l}{l!}}\right)^{\alpha}c^{xt} & =\sum\limits_{n=0}^{\infty}\frac{\partial^p}{\partial x^p}\left[\mathcal{Q}_{n}^{[m-1,\alpha]}(x,b,c;\lambda;u,v)\right]\frac{t^n}{n!}\\
\intertext{usando el \textit{Teorema \ref{par}} se sigue que}
t^p(\log c)^p\left(\frac{(2^ut^v)^{m}}{\lambda b^t+\sum\limits_{l=0}^{m-1}\frac{(t\log b)^l}{l!}}\right)^{\alpha}c^{xt} & =\sum\limits_{n=0}^{\infty}(\log c)^p\frac{n!}{(n-p)!}\mathcal{Q}_{n}^{[m-1,\alpha]}(x,b,c;\lambda;u,v)\frac{t^n}{n!}\\
\intertext{como $c\in \RR^{+}$ podemos multiplicar por $(\log c)^{-p}$ en ambos miembros y tenemos el final de la prueba:}
t^p\left(\frac{(2^ut^v)^{m}}{\lambda b^t+\sum\limits_{l=0}^{m-1}\frac{(t\log b)^l}{l!}}\right)^{\alpha}c^{xt} & =\sum\limits_{n=0}^{\infty}\frac{n!}{(n-p)!}\mathcal{Q}_{n-p}^{[m-1,\alpha]}(x,b,c;\lambda;u,v)\frac{t^n}{n!}.
\end{align*}
\end{proof}

\vspace{0.25cm}

\begin{lemma}
La nueva clase de polinomios tipo Apostol generalizados $\mathcal{Q}_n^{[m-1,\alpha]}(x,b,c;\lambda;u,v)$ satisface la siguiente relaci�n:
\begin{align}\label{e3}
& \nonumber\lambda \mathcal{Q}_n^{[m-1,\alpha]}(x+1,b,c;\lambda;u,v)+\mathcal{Q}_n^{[m-1,\alpha]}(x,b,c;\lambda;u,v)\\
& =\sum\limits_{k=0}^{n}\binom{n}{k}\binom{k}{v}v!(\log c)^v\mathcal{Q}_{k-v}^{[m-1,\alpha]}(x,b,c;\lambda;u,v)\mathcal{Q}_{n-k}^{(-1)}(0;\lambda;1,c,a;u,v),
\end{align}
donde $\mathcal{Q}_{n-k}^{(-1)}(0;\lambda;1,c,a;u,v)$ son los polinomios tipo Apostol generalizados definidos por (\ref{eqnewpol2}).
\end{lemma}

\vspace{0.25cm}

\begin{proof}
Por (\ref{newpol}) y el \textit{Teorema \ref{par}} cuando $p=1$, tenemos

\begin{align*}
& \displaystyle \sum\limits_{n=0}^{\infty}\left[\lambda \mathcal{Q}_n^{[m-1,\alpha]}(x+1,b,c;\lambda;u,v)+\mathcal{Q}_n^{[m-1,\alpha]}(x,b,c;\lambda;u,v)\right]\frac{t^n}{n!}\\
& = \displaystyle\lambda \sum\limits_{n=0}^{\infty}\mathcal{Q}_n^{[m-1,\alpha]}(x+1,b,c;\lambda;u,v)\frac{t^n}{n!}+\sum\limits_{n=0}^{\infty}\mathcal{Q}_n^{[m-1,\alpha]}(x,b,c;\lambda;u,v)\frac{t^n}{n!}\\
& =\lambda \displaystyle\left(\frac{(2^ut^v)^{m}}{\lambda b^t+\sum\limits_{l=0}^{m-1}\frac{(t\log b)^l}{l!}}\right)^{\alpha}c^{(x+1)t}+ \displaystyle\left(\frac{(2^ut^v)^{m}}{\lambda b^t+\sum\limits_{l=0}^{m-1}\frac{(t\log b)^l}{l!}}\right)^{\alpha}c^{xt}\\ 
& = \displaystyle\left(\frac{(2^ut^v)^{m}}{\lambda b^t+\sum\limits_{l=0}^{m-1}\frac{(t\log b)^l}{l!}}\right)^{\alpha}c^{xt}(\lambda c^t+1)=\left[2^ut^v\displaystyle\left(\frac{(2^ut^v)^{m}}{\lambda b^t+\sum\limits_{l=0}^{m-1}\frac{(t\log b)^l}{l!}}\right)^{\alpha}c^{xt}\right]\left(\displaystyle \frac{2^ut^v}{\lambda c^t+1}\right)^{(-1)}\\ 
& = \displaystyle 2^u\frac{\partial^v}{\partial x^v}\left[\left(\frac{(2^ut^v)^{m}}{\lambda b^t+\sum\limits_{l=0}^{m-1}\frac{(t\log b)^l}{l!}}\right)^{\alpha}c^{xt}\right] \left(\frac{2^ut^v}{\lambda c^t+1}\right)^{(-1)}c^{0t}\\
&= \displaystyle \frac{\partial^v}{\partial x^v}\left[ \sum\limits_{n=0}^{\infty}\mathcal{Q}_n^{[m-1,\alpha]}(x,b,c;\lambda;u,v)\frac{t^n}{n!}\right]\left(\sum\limits_{n=0}^{\infty}\mathcal{Q}_n^{(-1)}(0;\lambda;1,c,a;u,v)\frac{t^n}{n!}\right)\\
&=\displaystyle \left[ \sum\limits_{n=0}^{\infty}\frac{\partial^v}{\partial x^v}\left(\mathcal{Q}_n^{[m-1,\alpha]}(x,b,c;\lambda;u,v)\right)\frac{t^n}{n!}\right]\left(\sum\limits_{n=0}^{\infty}\mathcal{Q}_n^{(-1)}(0;\lambda;1,c,a;u,v)\frac{t^n}{n!}\right)\\ 
& = \displaystyle \left[ \sum\limits_{n=0}^{\infty}\frac{n!}{(n-v)!}(\log c)^v\mathcal{Q}_{n-v}^{[m-1,\alpha]}(x,b,c;\lambda;u,v)\frac{t^n}{n!}\right]\left(\sum\limits_{n=0}^{\infty}\mathcal{Q}_n^{(-1)}(0;\lambda;1,c,a;u,v)\frac{t^n}{n!}\right)\\
& =\displaystyle \sum\limits_{n=0}^{\infty} \sum\limits_{k=0}^{n}\binom{n}{k}\binom{k}{v}v!(\log c)^v\mathcal{Q}_{k-v}^{[m-1,\alpha]}(x,b,c;\lambda;u,v)\mathcal{Q}_{n-k}^{(-1)}(0;\lambda;1,c,a;u,v)\frac{t^n}{n!}.
\end{align*}
Comparando los coeficientes de $\frac{t^n}{n!}$ en ambos miembros de la ecuaci�n anterior, podemos obtener la identidad (\ref{e3}).
\end{proof}

\vspace{0.25cm}

\begin{teo}
La relaci�n
{\small
\begin{align}
& \nonumber\mathcal{Q}_n^{[m-1,\alpha]}(x+y,b,c;\lambda;u,v)\\
& \label{eqls}=\displaystyle \frac{1}{2}\sum\limits_{j=0}^{n}\binom{n}{j}\left[\lambda \mathcal{Q}_{n}^{[m-1,\alpha]}(y+1,b,c;\lambda;u,v)+(\log b)^{j}\mathcal{Q}_{n}^{[m-1,\alpha]}(y,b,c;\lambda;u,v)\right]\mathfrak{E}_{n-j}(x;\lambda).
\end{align}
}
se mantiene entre los polinomios $\mathcal{Q}_n^{[m-1,\alpha]}(x,b,c;\lambda;u,v)$ y los polinomios de Apostol-Euler $\mathfrak{E}_{n}(x;\lambda)$ definidos por (\ref{eqae})(con $\alpha = 1$).
\end{teo}

\begin{proof}
En primer lugar, sustituimos (\ref{lemeu}) en el lado derecho de la (\ref{eq9}) y obtenemos
\begin{align}
& \nonumber\mathcal{Q}_n^{[m-1,\alpha]}(x+y,b,c;\lambda;u,v)\\
&\nonumber =\displaystyle \sum\limits_{k=0}^{n}\binom{n}{k}\mathcal{Q}_k^{[m-1,\alpha]}(y,b,c;\lambda;u,v)(\log b)^{n-k}\frac{1}{2}\left[\lambda \sum\limits_{j=0}^{n-k}\binom{n-k}{j}\mathfrak{E}_{j}(x;\lambda)+\mathfrak{E}_{n-k}(x;\lambda)\right]\\
& \nonumber=\displaystyle \sum\limits_{k=0}^{n}\binom{n}{k}\mathcal{Q}_k^{[m-1,\alpha]}(y,b,c;\lambda;u,v)(\log b)^{n-k}\frac{1}{2}\lambda \sum\limits_{j=0}^{n-k}\binom{n-k}{j}\mathfrak{E}_{j}(x;\lambda)\\
& \label{eq89} + \displaystyle\sum\limits_{k=0}^{n}\binom{n}{k}\mathcal{Q}_k^{[m-1,\alpha]}(y,b,c;\lambda;u,v)(\log b)^{n-k}\frac{1}{2}\mathfrak{E}_{n-k}(x;\lambda).
\end{align}
La primera suma en (\ref{eq89}) es igual a

\begin{align}
&\nonumber\displaystyle \sum\limits_{k=0}^{n}\binom{n}{k}\mathcal{Q}_k^{[m-1,\alpha]}(y,b,c;\lambda;u,v)(\log b)^{n-k}\frac{1}{2}\lambda \sum\limits_{j=0}^{n-k}\binom{n-k}{j}\mathfrak{E}_{j}(x;\lambda) \\
&=\nonumber\displaystyle \sum\limits_{k=0}^{n}\sum\limits_{j=0}^{n-k}\binom{n}{k}(\log b)^{n-k}\frac{1}{2}\lambda \binom{n-k}{j}\mathcal{Q}_k^{[m-1,\alpha]}(y,b,c;\lambda;u,v)\mathfrak{E}_{j}(x;\lambda)\\
&=\nonumber\displaystyle \sum\limits_{j=0}^{n}\sum\limits_{k=0}^{n-j}\binom{n}{n-k}\binom{n-k}{j}(\log b)^{n-k}\frac{\lambda}{2}{F}_k^{[m-1,\alpha]}(y,b,c;\lambda;u,v)\mathfrak{E}_{j}(x;\lambda)\\
&=\nonumber\displaystyle \sum\limits_{j=0}^{n}\sum\limits_{k=0}^{n-j}\binom{n}{n-k}\binom{n-j}{n-j-k}(\log b)^{n-k}\frac{\lambda}{2}\mathcal{Q}_k^{[m-1,\alpha]}(y,b,c;\lambda;u,v)\mathfrak{E}_{j}(x;\lambda)
\end{align}

\begin{align}
&=\nonumber\displaystyle \sum\limits_{j=0}^{n}\frac{\lambda}{2}\binom{n}{j}\mathfrak{E}_{j}(x;\lambda)\sum\limits_{k=0}^{n-j}\binom{n-j}{k}\mathcal{Q}_k^{[m-1,\alpha]}(y,b,c;\lambda;u,v)(\log b)^{n-k}\\
&=\label{eq90}\displaystyle \sum\limits_{j=0}^{n}\frac{\lambda}{2}\binom{n}{j}\mathfrak{E}_{j}(x;\lambda)\mathcal{Q}_{n-j}^{[m-1,\alpha]}(y+1,b,c;\lambda;u,v).
\end{align}
Esto se logra invirtiendo el orden de la sumatoria, usando la ecuaci�n (\ref{eq9}) (cuando $x\longrightarrow 1; n\longrightarrow n-j$) y la aplicaci�n de las siguientes identidades elementales de combinatorios:
\begin{equation}
\binom{m}{l}\binom{l}{m}=\binom{m}{n}\binom{m-n}{l-n}\,\,\,\,\,(m\geq l\geq n: l,m,n \in \NN_{0}),
\end{equation}
\begin{equation}\label{comb}
\binom{n}{n-k}=\binom{n}{k}\,\,\,\,\,(n\geq k: k,n \in \NN_{0}),
\end{equation}
la segunda suma en (\ref{eq89}) es igual a
\begin{align}
& \nonumber\displaystyle \sum\limits_{k=0}^{n}\binom{n}{k}\mathcal{Q}_k^{[m-1,\alpha]}(y,b,c;\lambda;u,v)(\log b)^{n-k}\frac{1}{2}\mathfrak{E}_{n-k}(x;\lambda)\\
& \label{eq91}\displaystyle=\sum\limits_{k=0}^{n}\binom{n}{k}\mathcal{Q}_{n-k}^{[m-1,\alpha]}(y,b,c;\lambda;u,v)(\log b)^{k}\frac{1}{2}\mathfrak{E}_{k}(x;\lambda).
\end{align}
Combinando (\ref{eq90}) y (\ref{eq91})
\begin{align*}
& \mathcal{Q}_n^{[m-1,\alpha]}(x+y,b,c;\lambda;u,v)\\ 
& =\displaystyle \sum\limits_{j=0}^{n}\frac{\lambda}{2}\binom{n}{j}\mathfrak{E}_{j}(x;\lambda)\mathcal{Q}_{n-j}^{[m-1,\alpha]}(y+1,b,c;\lambda;u,v)\\
& \displaystyle +\sum\limits_{j=0}^{n}\binom{n}{j}\mathcal{Q}_{n-j}^{[m-1,\alpha]}(y,b,c;\lambda;u,v)(\log b)^{j}\frac{1}{2}\mathfrak{E}_{j}(x;\lambda)\\
& =\displaystyle \frac{1}{2}\sum\limits_{j=0}^{n}\binom{n}{j}\left[\lambda \mathcal{Q}_{n-j}^{[m-1,\alpha]}(y+1,b,c;\lambda;u,v)+(\log b)^{j}\mathcal{Q}_{n-j}^{[m-1,\alpha]}(y,b,c;\lambda;u,v)\right]\mathfrak{E}_{j}(x;\lambda)\\
& =\displaystyle \frac{1}{2}\sum\limits_{j=0}^{n}\binom{n}{j}\left[\lambda \mathcal{Q}_{n}^{[m-1,\alpha]}(y+1,b,c;\lambda;u,v)+(\log b)^{j}\mathcal{Q}_{n}^{[m-1,\alpha]}(y,b,c;\lambda;u,v)\right]\mathfrak{E}_{n-j}(x;\lambda).
\end{align*}
\end{proof}



\vspace{0.5cm}
\section{F�rmula de recurrencia y ecuaci�n diferencial}\label{3d2}
En esta secci�n ilustraremos la f�rmula de recurrencia y ecuaci�n diferencial que satisface $\mathcal{Q}_n^{[m-1,\alpha]}(x,b,c;\lambda;u,v)$ usando las identidades fundamentales de las series de potencias y las definiciones de polinomios de Appell.

\subsection{Ecuaci�n diferencial para $\mathcal{Q}_n^{[m-1,\alpha]}(x,b,c;\lambda;u,v)$}
\begin{teo}\label{teoeqdif}
Los polinomios $\mathcal{Q}_n^{[m-1,\alpha]}(x,b;\lambda;u,0)$ satisfacen la ecuaci�n diferencial:
\begin{equation}\label{eqappell}
\displaystyle \frac{\alpha_{n-1}}{(n-1)!}y^{(n)}+\frac{\alpha_{n-2}}{(n-2)!}y^{(n-1)}+\cdots+\frac{\alpha_{1}}{1!}y^{''}+(x+\alpha_0)y^{'}-ny=0,
\end{equation}
donde los coeficientes num�ricos $\alpha_{k}$, $k = 1,2,\cdots,n-1$ est�n definidos por
{\footnotesize
\begin{align}
\displaystyle \alpha_n& \nonumber =\frac{-\alpha(\log b)\lambda}{2^m}\frac{ n!}{(n+m)!}\mathcal{G}_{n+m}^{[m-1]}(1,b,b;\lambda)-\frac{\alpha(\log b)}{2^m}\left(\sum\limits_{k=0}^{m-1}\frac{(\log b)^k}{k!}\frac{n!}{(n+m-k)!}\mathcal{G}_{n+m-k}^{[m-1]}(b;\lambda)\right)\\
& +\frac{\alpha(\log b)^mn!}{2^m(m-1)!(n+1)!}\mathcal{G}_{n+1}^{[m-1]}(b;\lambda)
\end{align}
}
\end{teo}

\vspace{0.25cm}

\begin{proof}
Consideremos $\mathcal{Q}_n^{[m-1,\alpha]}(x,b,c;\lambda;u,v)$ y definamos
\begin{equation*}
\mathcal{Q}_n^{[m-1,\alpha]}(x,b;\lambda;u,v):=\mathcal{Q}_n^{[m-1,\alpha]}(x,b,e;\lambda;u,v)
\end{equation*}
\begin{equation*}
\mathcal{G}_n^{[m-1,\alpha]}(x,b;\lambda):=\mathcal{Q}_n^{[m-1,\alpha]}(x,b;\lambda;1,1).
\end{equation*}
Por otro lado tenemos que $\mathcal{Q}_n^{[m-1,\alpha]}(x,b;\lambda;u,v)$ es un polinomio de Appell ya que verifica (\ref{ruler}) como se mostr� en (\ref{e7}), por lo tanto satisface la ecuaci�n diferencial (\ref{eqappell}). A continuaci�n determinaremos las contantes $\alpha_k$, identificando $A(t)$ tenemos:
{\footnotesize
\begin{align*}
& A(t)=\displaystyle\left(\frac{(2^ut^v)^{m}}{\lambda b^t+\sum\limits_{l=0}^{m-1}\frac{(t\log b)^l}{l!}}\right)^{\alpha}\\
\intertext{{\normalsize derivando en ambos miembros, obtenemos}}
& A'(t)=\displaystyle\left[\frac{mv2^{mu}t^{mv-1}\left(\lambda b^t+\sum\limits_{l=0}^{m-1}\frac{(t\log b)^l}{l!}\right)-2^{mu}t^{mv}(\log b)\left(\lambda b^t+\sum\limits_{l=0}^{m-2}\frac{(t\log b)^{l}}{l!}\right)}{\left(\lambda b^t+\sum\limits_{l=0}^{m-1}\frac{(t\log b)^l}{l!}\right)^2}\right]\alpha\left(\frac{(2^ut^v)^{m}}{\lambda b^t+\sum\limits_{l=0}^{m-1}\frac{(t\log b)^l}{l!}}\right)^{\alpha-1}\\
& A'(t)=\displaystyle\left[\frac{mvt^{-1}\left(\lambda b^t+\sum\limits_{l=0}^{m-1}\frac{(t\log b)^l}{l!}\right)-(\log b)\left(\lambda b^t+\sum\limits_{l=0}^{m-2}\frac{(t\log b)^{l}}{l!}\right)}{\left(\lambda b^t+\sum\limits_{l=0}^{m-1}\frac{(t\log b)^l}{l!}\right)}\right]\alpha\left(\frac{(2^ut^v)^{m}}{\lambda b^t+\sum\limits_{l=0}^{m-1}\frac{(t\log b)^l}{l!}}\right)^{\alpha}\\
\intertext{{\normalsize dividiendo por $A(t)$, tenemos}}
& \displaystyle \frac{A'(t)}{A(t)} =\alpha\left[\frac{mvt^{-1}\left(\lambda b^t+\sum\limits_{l=0}^{m-1}\frac{(t\log b)^l}{l!}\right)-(\log b)\left(\lambda b^t+\sum\limits_{l=0}^{m-2}\frac{(t\log b)^{l}}{l!}\right)}{\left(\lambda b^t+\sum\limits_{l=0}^{m-1}\frac{(t\log b)^l}{l!}\right)}\right]\\
\intertext{{\normalsize haciendo $v=0$, tenemos}}
& \displaystyle \frac{A'(t)}{A(t)} =-\alpha\left[\frac{(\log b)\left(\lambda b^t+\sum\limits_{l=0}^{m-2}\frac{(t\log b)^{l}}{l!}\right)}{\left(\lambda b^t+\sum\limits_{l=0}^{m-1}\frac{(t\log b)^l}{l!}\right)}\right]\\
\intertext{{\normalsize separando la adici�n en el numerador del miembro de la derecha tenemos}}
& \displaystyle \frac{A'(t)}{A(t)} =-\alpha\left[\frac{(\log b)\lambda b^t}{\left(\lambda b^t+\sum\limits_{l=0}^{m-1}\frac{(t\log b)^l}{l!}\right)}\right]-\alpha\left[\frac{(\log b)\left(\sum\limits_{l=0}^{m-2}\frac{(t\log b)^{l}}{l!}\right)}{\left(\lambda b^t+\sum\limits_{l=0}^{m-1}\frac{(t\log b)^l}{l!}\right)}\right]\\
& \displaystyle \frac{A'(t)}{A(t)} =\frac{-\alpha(\log b)\lambda}{2^m}t^{-m}\left[\frac{ 2^mt^mb^t}{\left(\lambda b^t+\sum\limits_{l=0}^{m-1}\frac{(t\log b)^l}{l!}\right)}\right]-\alpha(\log b)\left[\sum\limits_{k=0}^{m-2}\frac{(t\log b)^k}{k!\left(\lambda b^t+\sum\limits_{l=0}^{m-1}\frac{(t\log b)^l}{l!}\right)}\right]\\
& \displaystyle \frac{A'(t)}{A(t)} =\frac{-\alpha(\log b)\lambda}{2^m}t^{-m}\left[\frac{ 2^mt^mb^t}{\left(\lambda b^t+\sum\limits_{l=0}^{m-1}\frac{(t\log b)^l}{l!}\right)}\right]-\frac{\alpha(\log b)}{2^m}\left[\sum\limits_{k=0}^{m-2}\frac{(\log b)^k}{k!}t^{k-m}\frac{2^mt^m}{\left(\lambda b^t+\sum\limits_{l=0}^{m-1}\frac{(t\log b)^l}{l!}\right)}\right]\\
\intertext{{\normalsize usando el \textit{Teorema \ref{par2}}}, tenemos}
& \displaystyle \frac{A'(t)}{A(t)}=\frac{-\alpha(\log b)\lambda}{2^m}\sum\limits_{n=0}^{\infty}\frac{ n!}{(n+m)!}\mathcal{G}_{n+m}^{[m-1]}(1,b,b;\lambda)\frac{t^n}{n!}-\frac{\alpha(\log b)}{2^m}\left(\sum\limits_{k=0}^{m-2}\frac{(\log b)^k}{k!}\sum\limits_{n=0}^{\infty}\frac{ n!}{(n+m-k)!}\mathcal{G}_{n+m-k}^{[m-1]}(b;\lambda)\frac{t^n}{n!}\right)\\
& \displaystyle \frac{A'(t)}{A(t)}=\displaystyle \sum\limits_{n=0}^{\infty}\left[\frac{-\alpha(\log b)\lambda}{2^m}\frac{ n!}{(n+m)!}\mathcal{G}_{n+m}^{[m-1]}(1,b,b;\lambda)-\frac{\alpha(\log b)}{2^m}\left(\sum\limits_{k=0}^{m-2}\frac{(\log b)^k}{k!}\frac{n!}{(n+m-k)!}\mathcal{G}_{n+m-k}^{[m-1]}(b;\lambda)\right)\right]\frac{t^n}{n!}\\
\intertext{{\normalsize usando (\ref{eqalpha}), tenemos que}}
\displaystyle \alpha_n& =\frac{-\alpha(\log b)\lambda}{2^m}\frac{ n!}{(n+m)!}\mathcal{G}_{n+m}^{[m-1]}(1,b,b;\lambda)-\frac{\alpha(\log b)}{2^m}\left(\sum\limits_{k=0}^{m-2}\frac{(\log b)^k}{k!}\frac{n!}{(n+m-k)!}\mathcal{G}_{n+m-k}^{[m-1]}(b;\lambda)\right)\\
\intertext{{\normalsize finalmente}}
\displaystyle \alpha_n& =\frac{-\alpha(\log b)\lambda}{2^m}\frac{ n!}{(n+m)!}\mathcal{G}_{n+m}^{[m-1]}(1,b,b;\lambda)-\frac{\alpha(\log b)}{2^m}\left(\sum\limits_{k=0}^{m-1}\frac{(\log b)^k}{k!}\frac{n!}{(n+m-k)!}\mathcal{G}_{n+m-k}^{[m-1]}(b;\lambda)\right)\\
& +\frac{\alpha(\log b)}{2^m}\frac{(\log b)^{m-1}}{(m-1)!}\frac{n!}{(n+1)!}\mathcal{G}_{n+1}^{[m-1]}(b;\lambda)
\end{align*}
}
\end{proof}


\vspace{0.25cm}

Especialmente haciendo $u=1$ en el \textit{Teorema \ref{teoeqdif}} tenemos el siguiente corolario.

\vspace{0.25cm}

\begin{coro}\label{eqdif2}
Los polinomios $\mathcal{E}_n^{[m-1,\alpha]}(x,b;\lambda)$ satisfacen la ecuaci�n diferencial:
\begin{equation}\label{eqappell2}
\displaystyle \frac{\alpha_{n-1}}{(n-1)!}y^{(n)}+\frac{\alpha_{n-2}}{(n-2)!}y^{(n-1)}+\cdots+\frac{\alpha_{1}}{1!}y^{''}+(x+\alpha_0)y^{'}-ny=0,
\end{equation}
donde los coeficientes num�ricos $\alpha_{k}$, $k = 1,2,\cdots,n-1$ est�n definidos por
{\footnotesize
\begin{align}
\displaystyle \alpha_n& \nonumber =\frac{-\alpha(\log b)\lambda}{2^m}\frac{ n!}{(n+m)!}\mathcal{G}_{n+m}^{[m-1]}(1,b,b;\lambda)-\frac{\alpha(\log b)}{2^m}\left(\sum\limits_{k=0}^{m-1}\frac{(\log b)^k}{k!}\frac{n!}{(n+m-k)!}\mathcal{G}_{n+m-k}^{[m-1]}(b;\lambda)\right)\\
& +\frac{\alpha(\log b)^mn!}{2^m(m-1)!(n+1)!}\mathcal{G}_{n+1}^{[m-1]}(b;\lambda)
\end{align}
}
\end{coro}

\vspace{0.25cm}

\begin{obs}
Si hacemos $m=1$ y $b=e$ en el \textit{Corolario \ref{eqdif2}} tenemos que los polinomios $\mathcal{E}_n^{(\alpha)}(x;\lambda)$ satisfacen la ecuaci�n diferencial:
\begin{equation}\label{eqappell3}
\displaystyle \frac{\alpha_{n-1}}{(n-1)!}y^{(n)}+\frac{\alpha_{n-2}}{(n-2)!}y^{(n-1)}+\cdots+\frac{\alpha_{1}}{1!}y^{''}+(x+\alpha_0)y^{'}-ny=0,
\end{equation}
donde los coeficientes num�ricos $\alpha_{k}$, $k = 1,2,\cdots,n-1$ est�n definidos por
\begin{align}
\displaystyle \alpha_n&  =-\frac{\alpha\lambda\mathcal{G}_{n+1}(1;\lambda)}{2(n+1)} \nonumber
\end{align}
la cual es justamente el \textit{Corolario 2.8} de \cite{atp1}
\end{obs}


\vspace{0.25cm}

Haciendo $m=1$ y $b=e$ en el \textit{Teorema \ref{teoeqdif}} tenemos el siguiente corolario.

\vspace{0.25cm}

\begin{coro}\label{teoeqdif1}
Los polinomios $\mathcal{F}_n^{(\alpha)}(x;\lambda;u,0)$ satisfacen la ecuaci�n diferencial:
\begin{equation}\label{eqappell1}
\displaystyle \frac{\alpha_{n-1}}{(n-1)!}y^{(n)}+\frac{\alpha_{n-2}}{(n-2)!}y^{(n-1)}+\cdots+\frac{\alpha_{1}}{1!}y^{''}+(x+\alpha_0)y^{'}-ny=0,
\end{equation}
donde los coeficientes num�ricos $\alpha_{k}$, $k = 1,2,\cdots,n-1$ est�n definidos por
\begin{align}
\displaystyle \alpha_n&  =-\frac{\alpha\lambda\mathcal{G}_{n+1}(1;\lambda)}{2(n+1)}
\intertext{que es justamente el \textit{Teorema 2.7} de \cite{atp1}} \nonumber
\end{align}
\end{coro}

\vspace{0.25cm}

\subsection{F�rmula de recurrencia para $\mathcal{Q}_n^{[m-1,\alpha]}(x,b,c;\lambda;u,v)$}
\begin{teo}
La nueva clase de polinomios tipo Apostol generalizados verifica la siguiente relaci�n de recurrencia
\begin{align}
& \nonumber \left(\log \frac{b^{\alpha}}{c^x} \right)\mathcal{Q}_{n}^{[m-1,\alpha]}(x,b,c;\lambda;u,v)=\left(\frac{\alpha mv}{n+1}-1\right)\mathcal{Q}_{n+1}^{[m-1,\alpha]}(x,b,c;\lambda;u,v)\\
& \label{relrec} +\frac{\alpha (\log b)^m}{2^{um}(m-1)!}\frac{n!}{(n+vm-m+1)!}\mathcal{Q}_{n+vm-m+1}^{[m-1,\alpha+1]}(x,b,c;\lambda;u,v)
\end{align}
\end{teo}


\vspace{0.25cm}

\begin{proof}
Consideremos 
{\footnotesize
\begin{align*}
& \displaystyle\left(\frac{(2^ut^v)^{m}}{\lambda b^t+\sum\limits_{l=0}^{m-1}\frac{(t\log b)^l}{l!}}\right)^{\alpha}c^{xt} =\displaystyle\sum\limits_{n=0}^{\infty}
\mathcal{Q}_n^{[m-1,\alpha]}(x,b,c;\lambda;u,v)\frac{t^n}{n!}\\
\intertext{{\normalsize derivando respecto a $t$ en ambos miembros, tenemos}}
& \displaystyle\frac{\partial}{\partial t}\left[\left(\frac{(2^ut^v)^{m}}{\lambda b^t+\sum\limits_{l=0}^{m-1}\frac{(t\log b)^l}{l!}}\right)^{\alpha}\right]c^{xt} +\left(\frac{(2^ut^v)^{m}}{\lambda b^t+\sum\limits_{l=0}^{m-1}\frac{(t\log b)^l}{l!}}\right)^{\alpha}\frac{\partial}{\partial t}\left[c^{xt}\right] =\displaystyle\frac{\partial}{\partial t}\left[\sum\limits_{n=0}^{\infty}
\mathcal{Q}_n^{[m-1,\alpha]}(x,b,c;\lambda;u,v)\frac{t^n}{n!}\right]\\
\intertext{{\normalsize resolviendo las derivadas ambos miembros, se sigue que}}
& \displaystyle\left[\frac{mv2^{mu}t^{mv-1}\left(\lambda b^t+\sum\limits_{l=0}^{m-1}\frac{(t\log b)^l}{l!}\right)-2^{mu}t^{mv}(\log b)\left(\lambda b^t+\sum\limits_{l=0}^{m-2}\frac{(t\log b)^{l}}{l!}\right)}{\left(\lambda b^t+\sum\limits_{l=0}^{m-1}\frac{(t\log b)^l}{l!}\right)^2}\right]c^{xt}\alpha\left(\frac{(2^ut^v)^{m}}{\lambda b^t+\sum\limits_{l=0}^{m-1}\frac{(t\log b)^l}{l!}}\right)^{\alpha-1}\\
& +\displaystyle \left(\frac{(2^ut^v)^{m}}{\lambda b^t+\sum\limits_{l=0}^{m-1}\frac{(t\log b)^l}{l!}}\right)^{\alpha}xc^{xt}(\log c)=\sum\limits_{n=0}^{\infty}
\mathcal{Q}_{n+1}^{[m-1,\alpha]}(x,b,c;\lambda;u,v)\frac{t^n}{n!}\\
\intertext{{\normalsize resolviendo la primera suma y expresando como serie de potencia la segunda suma, tenemos}}
& \displaystyle\left[\frac{mvt^{-1}\left(\lambda b^t+\sum\limits_{l=0}^{m-1}\frac{(t\log b)^l}{l!}\right)-(\log b)\left(\lambda b^t+\sum\limits_{l=0}^{m-2}\frac{(t\log b)^{l}}{l!}\right)}{(2^ut^v)^{m}}\right]c^{xt}\alpha\left(\frac{(2^ut^v)^{m}}{\lambda b^t+\sum\limits_{l=0}^{m-1}\frac{(t\log b)^l}{l!}}\right)^{\alpha+1}\\
& +\displaystyle \sum\limits_{n=0}^{\infty}
x(\log c)\mathcal{Q}_{n}^{[m-1,\alpha]}(x,b,c;\lambda;u,v)\frac{t^n}{n!}=\sum\limits_{n=0}^{\infty}
\mathcal{Q}_{n+1}^{[m-1,\alpha]}(x,b,c;\lambda;u,v)\frac{t^n}{n!}\\
\intertext{{\normalsize expresando la primera suma en t�rminos de $m-1$, se sigue que}}
& \displaystyle\left[\frac{mvt^{-1}\left(\lambda b^t+\sum\limits_{l=0}^{m-1}\frac{(t\log b)^l}{l!}\right)-(\log b)\left(\lambda b^t+\sum\limits_{l=0}^{m-1}\frac{(t\log b)^{l}}{l!}\right)+(\log b)\left(\frac{(t\log b)^{m-1}}{(m-1)!}\right)}{(2^ut^v)^{m}}\right]c^{xt}\alpha\left(\frac{(2^ut^v)^{m}}{\lambda b^t+\sum\limits_{l=0}^{m-1}\frac{(t\log b)^l}{l!}}\right)^{\alpha+1}\\
& +\displaystyle \sum\limits_{n=0}^{\infty}
x(\log c)\mathcal{Q}_{n}^{[m-1,\alpha]}(x,b,c;\lambda;u,v)\frac{t^n}{n!}=\sum\limits_{n=0}^{\infty}
\mathcal{Q}_{n+1}^{[m-1,\alpha]}(x,b,c;\lambda;u,v)\frac{t^n}{n!}\\
\intertext{{\normalsize ahora, por aplicaci�n del \textit{Teorema \ref{par2}} y expresando la primera suma como serie de potencia, tenemos}}
& \displaystyle \sum\limits_{n=0}^{\infty}\frac{\alpha mv}{n+1}\mathcal{Q}_{n+1}^{[m-1,\alpha]}(x,b,c;\lambda;u,v)\frac{t^n}{n!}-\sum\limits_{n=0}^{\infty}\alpha (\log b)\mathcal{Q}_{n}^{[m-1,\alpha]}(x,b,c;\lambda;u,v)\frac{t^n}{n!}\\
& +\sum\limits_{n=0}^{\infty}\frac{\alpha (\log b)^m}{2^{um}(m-1)!}\frac{n!}{(n+vm-m+1)!}\mathcal{Q}_{n+vm-m+1}^{[m-1,\alpha+1]}(x,b,c;\lambda;u,v)\frac{t^n}{n!}+\displaystyle \sum\limits_{n=0}^{\infty}
x(\log c)\mathcal{Q}_{n}^{[m-1,\alpha]}(x,b,c;\lambda;u,v)\frac{t^n}{n!}\\
& =\sum\limits_{n=0}^{\infty}
\mathcal{Q}_{n+1}^{[m-1,\alpha]}(x,b,c;\lambda;u,v)\frac{t^n}{n!}\\
\intertext{{\normalsize finalmente por identidad de series de potencias concluimos la prueba}}
& \left(\alpha \log b-x\log c \right)\mathcal{Q}_{n}^{[m-1,\alpha]}(x,b,c;\lambda;u,v)=\left(\frac{\alpha mv}{n+1}-1\right)\mathcal{Q}_{n+1}^{[m-1,\alpha]}(x,b,c;\lambda;u,v)\\
& +\frac{\alpha (\log b)^m}{2^{um}(m-1)!}\frac{n!}{(n+vm-m+1)!}\mathcal{Q}_{n+vm-m+1}^{[m-1,\alpha+1]}(x,b,c;\lambda;u,v)
\end{align*}
}
\end{proof}





\vspace{0.5cm}
\section{F�rmulas de conexi�n  entre $\mathcal{Q}_n^{[m-1,\alpha]}(x,b,c;\lambda;u,v)$ y otros polinomios y n�meros}\label{3e}
En esta secci�n se estudian algunas conexiones entre los polinomios $\mathcal{Q}_n^{[m-1,\alpha]}(x,b,c;\lambda;u,v)$ y otras familias de polinomios como los polinomios de Genocchi, los polinomios de Jacobi, los polinomios de Laguerre, los polinomios de Hermite, los polinomios de Bernoulli generalizados $B_{n}^{[m-1]}(x)$, los polinomios de Charlier y los polinomios de Bessel. Por �ltimo, se establece la relaci�n de la nueva clase de polinomios tipo Apostol generalizados con los n�meros de Stirling de segunda clase, todo ello siguiendo las mismas ideas usadas en la prueba de los teoremas de la secci�n anterior, y tambi�n se utilizan algunas de las siguientes propiedades b�sicas de las sumatorias:
\begin{equation}\label{s4}
\displaystyle \sum\limits_{j=0}^{n}A_{n-j}B_{n}=\sum\limits_{j=0}^{n}A_{j}B_{n-j},
\end{equation}
\begin{equation}\label{s3}
\displaystyle \sum\limits_{j=0}^{n-k}A_{j}=\sum\limits_{j=k}^{n}A_{n-j},
\end{equation}
\begin{equation}\label{s2}
\displaystyle \sum\limits_{j=0}^{n}\sum\limits_{k=0}^{n}a_{j}a_{k}=\left(\sum\limits_{j=0}^{n}a_{j}\right)\left(\sum\limits_{k=0}^{n}a_{k}\right) 
\end{equation}
y
\begin{equation}\label{s1}
\displaystyle \sum\limits_{j=0}^{n}\sum\limits_{k=0}^{n-j}a_{jk}=\sum\limits_{k=0}^{n}\sum\limits_{j=0}^{n-k}a_{jk}.\\
\end{equation}

\vspace{0.25cm}

\begin{teo}\label{teogen}
La relaci�n
{\footnotesize
\begin{align}\label{eq11}
& \nonumber\mathcal{Q}_n^{[m-1,\alpha]}(x+y,b,c;\lambda;u,v)\\
&=\displaystyle\frac{1}{2}\sum\limits_{k=0}^{n}
\frac{(\log c)^{k}}{k+1}\left[\binom{n}{k}\mathcal{Q}_{n-k}^{[m-1,\alpha]}(y,b,c;\lambda;u,v)+\sum\limits_{j=k}^{n}\binom{n}{j}\binom{j}{k}\mathcal{Q}_{n-j}^{[m-1,\alpha]}(y,b,c;\lambda;u,v)(\log c)^{j-k}\right]G_{k+1}(x).
\end{align}
}
se mantiene entre los polinomios $\mathcal{Q}_n^{[m-1,\alpha]}(x,b,c;\lambda;u,v)$ y los polinomios de Genocchi $G_{n}(x)$ definidos por (\ref{genpol23}). 
\end{teo}

\vspace{0.25cm}

\begin{proof} 
{\footnotesize
Por sustituci�n de (\ref{eq12}) en el lado derecho de (\ref{eq9}), tenemos
\begin{align*}
&\mathcal{Q}_n^{[m-1,\alpha]}(x+y,b,c;\lambda;u,v)\\ &=\displaystyle\sum
\limits_{j=0}^{n}
\binom{n}{j}\mathcal{Q}_j^{[m-1,\alpha]}(y,b,c;\lambda;u,v)\frac{(\log c)^{n-j}}{2(n-j+1)}\left[\sum\limits_{k=0}^{n-j}\binom{n-j+1}{k+1}G_{k+1}(x)+G_{n-j+1}(x)\right]\\
& =\displaystyle\sum
\limits_{j=0}^{n}
\binom{n}{j}\mathcal{Q}_j^{[m-1,\alpha]}(y,b,c;\lambda;u,v))\frac{(\log c)^{n-j}}{2(n-j+1)}\sum\limits_{k=0}^{n-j}\binom{n-j+1}{k+1}G_{k+1}(x)\\
&+ \displaystyle\sum
\limits_{j=0}^{n}
\binom{n}{j}\mathcal{Q}_j^{[m-1,\alpha]}(y,b,c;\lambda;u,v)\frac{(\log c)^{n-j}}{2(n-j+1)}G_{n-j+1}(x)\\
\intertext{usando (\ref{s2}) tenemos}
& =\displaystyle\sum
\limits_{j=0}^{n}\sum\limits_{k=0}^{n-j}\binom{n}{j}\mathcal{Q}_j^{[m-1,\alpha]}(y,b,c;\lambda;u,v)\frac{(\log c)^{n-j}}{2(n-j+1)}\binom{n-j+1}{k+1}G_{k+1}(x)\\
&+ \displaystyle\sum
\limits_{j=0}^{n}
\binom{n}{j}\mathcal{Q}_j^{[m-1,\alpha]}(y,b,c;\lambda;u,v)\frac{(\log c)^{n-j}}{2(n-j+1)}G_{n-j+1}(x)\\
\intertext{usando (\ref{s1}) en la primera suma y cambiando $j$ por $k$ se sigue que}
& =\displaystyle\sum
\limits_{k=0}^{n}\sum\limits_{j=0}^{n-k}\binom{n}{j}\mathcal{Q}_j^{[m-1,\alpha]}(y,b,c;\lambda;u,v)\frac{(\log c)^{n-j}}{2(n-j+1)}\binom{n-j+1}{k+1}G_{k+1}(x)\\
&+ \displaystyle\sum
\limits_{k=0}^{n}
\binom{n}{k}\mathcal{Q}_j^{[m-1,\alpha]}(y,b,c;\lambda;u,v))\frac{(\log c)^{n-k}}{2(n-k+1)}G_{n-k+1}(x)\\
\intertext{ahora, usando (\ref{s3}) en la primera suma y (\ref{s4}) en la segunda suma se sigue que}
& =\displaystyle\sum
\limits_{k=0}^{n}\sum\limits_{j=k}^{n}\binom{n}{n-j}\mathcal{Q}_j^{[m-1,\alpha]}(y,b,c;\lambda;u,v)\frac{(\log c)^{j}}{2(j+1)}\binom{j+1}{k+1}G_{k+1}(x)\\
&+ \displaystyle\sum
\limits_{k=0}^{n}
\binom{n}{n-k}\mathcal{Q}_{n-k}^{[m-1,\alpha]}(y,b,c;\lambda;u,v)\frac{(\log c)^{k}}{2(k+1)}G_{k+1}(x)\\
\intertext{amplificando por $\displaystyle \binom{j}{k}$ en la primera suma, tenemos}
& =\displaystyle\sum
\limits_{k=0}^{n}\sum\limits_{j=k}^{n}\binom{n}{j}\binom{j}{k}\mathcal{Q}_{n-j}^{[m-1,\alpha]}(y,b,c;\lambda;u,v)\frac{(\log c)^{j}}{2(j+1)}\binom{j+1}{k+1}\binom{j}{k}^{-1}G_{k+1}(x)\\
&+ \displaystyle\sum
\limits_{k=0}^{n}
\binom{n}{n-k}\mathcal{Q}_{n-k}^{[m-1,\alpha]}(y,b,c;\lambda;u,v)\frac{(\log c)^{k}}{2(k+1)}G_{k+1}(x)\\
\intertext{desarrolando el combinatorio en la primera suma}
& =\displaystyle\sum
\limits_{k=0}^{n}\sum\limits_{j=k}^{n}\binom{n}{j}\binom{j}{k}\mathcal{Q}_{n-j}^{[m-1,\alpha]}(y,b,c;\lambda;u,v)\frac{(\log c)^{j}}{2(k+1)}G_{k+1}(x)\\
&+ \displaystyle\sum
\limits_{k=0}^{n}
\binom{n}{k}\mathcal{Q}_{n-k}^{[m-1,\alpha]}(y,b,c;\lambda;u,v)\frac{(\log c)^{k}}{2(k+1)}G_{k+1}(x),\\
\intertext{finalmente, 
factorizando tenemos}
& =\displaystyle\frac{1}{2}\sum\limits_{k=0}^{n}
\frac{(\log c)^{k}}{k+1}\left[\sum\limits_{j=k}^{n}\binom{n}{j}\binom{j}{k}\mathcal{Q}_{n-j}^{[m-1,\alpha]}(y,b,c;\lambda;u,v)(\log c)^{j-k}+\binom{n}{k}\mathcal{Q}_{n-k}^{[m-1,\alpha]}(y,b,c;\lambda;u,v)\right]G_{k+1}(x).
\end{align*}
}
\end{proof}

\vspace{0.25cm}

Tomando $\lambda=:-\lambda, u=0$ y $v=1$ en el \textit{Teorema \ref{teogen}}, y luego multiplicando por $(-1)^\alpha$ en ambos lados del resultado, tenemos el siguiente corolario.

\vspace{0.25cm}

\begin{coro}\label{corber}
{\footnotesize
\begin{align}\label{eq121}
& \nonumber\mathcal{B}_n^{[m-1,\alpha]}(x+y,b,c;\lambda)\\
&=\displaystyle\frac{1}{2}\sum\limits_{k=0}^{n}
\frac{(\log c)^{k}}{k+1}\left[\binom{n}{k}\mathcal{B}_{n-k}^{[m-1,\alpha]}(y,b,c;\lambda)+\sum\limits_{j=k}^{n}\binom{n}{j}\binom{j}{k}\mathcal{B}_{n-j}^{[m-1,\alpha]}(y,b,c;\lambda)(\log c)^{j-k}\right]G_{k+1}(x).
\end{align}
}
\end{coro}

\vspace{0.25cm}

\begin{obs}
Si hacemos $b=c=e$ en el \textit{Corolario \ref{corber}} tenemos
{\footnotesize
\begin{align}\label{eq113}
& \nonumber\mathcal{B}_n^{[m-1,\alpha]}(x+y;\lambda)\\
&=\displaystyle\frac{1}{2}\sum\limits_{k=0}^{n}
\frac{1}{k+1}\left[\binom{n}{k}\mathcal{B}_{n-k}^{[m-1,\alpha]}(y;\lambda)+\sum\limits_{j=k}^{n}\binom{n}{j}\binom{j}{k}\mathcal{B}_{n-j}^{[m-1,\alpha]}(y;\lambda)\right]G_{k+1}(x).
\end{align}
}
\end{obs}

\vspace{0.25cm}

Tomando $u=1$ y $v=0$ en el \textit{Teorema \ref{teogen}}, tenemos el siguiente corolario.

\vspace{0.25cm}

\begin{coro}\label{coreul}
{\footnotesize
\begin{align}\label{eq1187}
& \nonumber\mathcal{E}_n^{[m-1,\alpha]}(x+y,b,c;\lambda)\\
&=\displaystyle\frac{1}{2}\sum\limits_{k=0}^{n}
\frac{(\log c)^{k}}{k+1}\left[\binom{n}{k}\mathcal{E}_{n-k}^{[m-1,\alpha]}(y,b,c;\lambda)+\sum\limits_{j=k}^{n}\binom{n}{j}\binom{j}{k}\mathcal{E}_{n-j}^{[m-1,\alpha]}(y,b,c;\lambda)(\log c)^{j-k}\right]G_{k+1}(x).
\end{align}
}
\end{coro}

\vspace{0.25cm}

\begin{obs}
Si hacemos $b=c=e$ en el \textit{Corolario \ref{coreul}} tenemos
{\footnotesize
\begin{align}\label{eq1123}
& \nonumber\mathfrak{E}_n^{[m-1,\alpha]}(x+y;\lambda)\\
&=\displaystyle\frac{1}{2}\sum\limits_{k=0}^{n}
\frac{1}{k+1}\left[\binom{n}{k}\mathfrak{E}_{n-k}^{[m-1,\alpha]}(y;\lambda)+\sum\limits_{j=k}^{n}\binom{n}{j}\binom{j}{k}\mathfrak{E}_{n-j}^{[m-1,\alpha]}(y;\lambda)\right]G_{k+1}(x).
\end{align}
} la cual es justamente la Ec.(4.1) de \cite{Ming}
\end{obs}

\vspace{0.25cm}

Al establecer $u=1$ y $v=1$ en el \textit{Teorema \ref{teogen}}, tenemos el siguiente corolario.

\vspace{0.25cm}

\begin{coro}\label{corgen}
{\footnotesize
\begin{align}\label{eq115}
& \nonumber\mathcal{G}_n^{[m-1,\alpha]}(x+y,b,c;\lambda)\\
&=\displaystyle\frac{1}{2}\sum\limits_{k=0}^{n}
\frac{(\log c)^{k}}{k+1}\left[\binom{n}{k}\mathcal{G}_{n-k}^{[m-1,\alpha]}(y,b,c;\lambda)+\sum\limits_{j=k}^{n}\binom{n}{j}\binom{j}{k}\mathcal{G}_{n-j}^{[m-1,\alpha]}(y,b,c;\lambda)(\log c)^{j-k}\right]G_{k+1}(x).
\end{align}
}
\end{coro}

\vspace{0.25cm}

\begin{obs}
Si hacemos $b=c=e$ en el \textit{Corolario \ref{corgen}} tenemos
{\footnotesize
\begin{align}\label{eq117}
& \nonumber\mathfrak{G}_n^{[m-1,\alpha]}(x+y;\lambda)\\
&=\displaystyle\frac{1}{2}\sum\limits_{k=0}^{n}
\frac{1}{k+1}\left[\binom{n}{k}\mathfrak{G}_{n-k}^{[m-1,\alpha]}(y;\lambda)+\sum\limits_{j=k}^{n}\binom{n}{j}\binom{j}{k}\mathfrak{G}_{n-j}^{[m-1,\alpha]}(y;\lambda)\right]G_{k+1}(x).
\end{align}
}
\end{obs}

\vspace{0.25cm}

\begin{teo}\label{teojac}
La relaci�n
{\footnotesize
\begin{align}
& \nonumber \mathcal{Q}_{n}^{[m-1,\mu]}(x+y,b,c;\lambda;u,v)\\
&\label{eq17} =\displaystyle\sum\limits_{k=0}^{n}(-1)^k
\sum\limits_{j=k}^{n}j!(\log c)^{j}\binom{j+\alpha}{j-k}\binom{n}{j}\frac{\alpha+\beta+2k+1}{(\alpha+\beta+k+1)_{j+1}}    \mathcal{Q}_{n-j}^{[m-1,\mu]}(y,b,c;\lambda;u,v)P_k^{(\alpha,\beta)}(1-2x)
\end{align}
}
se mantiene entre los polinomios $\mathcal{Q}_{n}^{[m-1,\mu]}(x,b,c;\lambda;u,v)$ y los polinomios de Jacobi $P_n^{(\alpha,\beta)}(x)$ definidos por (\ref{eq30}).
\end{teo}

\begin{proof}
{\footnotesize Por sustituci�n de (\ref{eq31}) en el lado derecho de (\ref{eq9}), tenemos
\begin{align*}
& \mathcal{Q}_{n}^{[m-1,\mu]}(x+y,b,c;\lambda;u,v)\\
& =\displaystyle\sum
\limits_{j=0}^{n}
\binom{n}{j}\mathcal{Q}_{j}^{[m-1,\mu]}(y,b,c;\lambda;u,v)(n-j)!(\log c)^{n-j}\displaystyle\sum\limits_{k=0}^{n-j}(-1)^k\binom{n-j+\alpha}{n-j-k}
\frac{\alpha+\beta+2k+1}{(\alpha+\beta+k+1)_{n-j+1}}P_k^{(\alpha,\beta)}(1-2x)\\
\intertext{aplicando (\ref{s2}), tenemos}
& =\displaystyle\sum
\limits_{j=0}^{n}\sum\limits_{k=0}^{n-j}
\binom{n}{j}\mathcal{Q}_{j}^{[m-1,\mu]}(y,b,c;\lambda;u,v)(n-j)!(\log c)^{n-j}(-1)^k\binom{n-j+\alpha}{n-j-k}
\frac{\alpha+\beta+2k+1}{(\alpha+\beta+k+1)_{n-j+1}}P_k^{(\alpha,\beta)}(1-2x)\\
\intertext{usando (\ref{s1}) se sigue que}
& =\displaystyle\sum
\limits_{k=0}^{n}\sum\limits_{j=0}^{n-k}
\binom{n}{j}\mathcal{Q}_{j}^{[m-1,\mu]}(y,b,c;\lambda;u,v)(n-j)!(\log c)^{n-j}(-1)^k\binom{n-j+\alpha}{n-j-k}
\frac{\alpha+\beta+2k+1}{(\alpha+\beta+k+1)_{n-j+1}}P_k^{(\alpha,\beta)}(1-2x)\\
\intertext{por uso de (\ref{s2}) nuevamente, obtenemos}
& =\displaystyle\sum
\limits_{k=0}^{n}(-1)^k\sum\limits_{j=0}^{n-k}
\binom{n}{j}\binom{n-j+\alpha}{n-j-k}\mathcal{Q}_{j}^{[m-1,\mu]}(y,b,c;\lambda;u,v)(n-j)!
(\log c)^{n-j}\frac{\alpha+\beta+2k+1}{(\alpha+\beta+k+1)_{n-j+1}}P_k^{(\alpha,\beta)}(1-2x)\\
\intertext{aplicando (\ref{s3}) se sigue que}
& =\displaystyle\sum
\limits_{k=0}^{n}(-1)^k\sum\limits_{j=0}^{n-k}
\binom{n}{n-j}\binom{j+\alpha}{j-k}\mathcal{Q}_{j}^{[m-1,\mu]}(y,b,c;\lambda;u,v)j!
(\log c)^{j}\frac{\alpha+\beta+2k+1}{(\alpha+\beta+k+1)_{n-j+1}}P_k^{(\alpha,\beta)}(1-2x)\\
\intertext{finalmenete, por (\ref{comb}) y ordenando los combinatorios, esto completa la prueba}
& =\displaystyle\sum\limits_{k=0}^{n}(-1)^k
\sum\limits_{j=k}^{n}j!(\log c)^{j}\binom{j+\alpha}{j-k}\binom{n}{j}\frac{\alpha+\beta+2k+1}{(\alpha+\beta+k+1)_{j+1}}    \mathcal{Q}_{n-j}^{[m-1,\mu]}(y,b,c;\lambda;u,v)P_k^{(\alpha,\beta)}(1-2x).
\end{align*}
}
\end{proof}

\vspace{0.25cm}

Al establecer $\lambda=:-\lambda, u=0$ y $v=1$ en el \textit{Teorema \ref{teojac}}, y luego multiplicar por $(-1)^\mu$ tenemos el siguiente corolario.

\vspace{0.25cm}

\begin{coro}\label{corber1}
{\footnotesize
\begin{align}
& \nonumber \mathcal{B}_{n}^{[m-1,\mu]}(x+y,b,c;\lambda)\\
&\nonumber =\displaystyle\sum\limits_{k=0}^{n}(-1)^k
\sum\limits_{j=k}^{n}j!(\log c)^{j}\binom{j+\alpha}{j-k}\binom{n}{j}\frac{\alpha+\beta+2k+1}{(\alpha+\beta+k+1)_{j+1}}    \mathcal{B}_{n-j}^{[m-1,\mu]}(y,b,c;\lambda)P_k^{(\alpha,\beta)}(1-2x)
\end{align}
}
\end{coro}

\vspace{0.25cm}

\begin{obs}
Si hacemos $b=c=e$ en el \textit{Corolario \ref{corber1}} tenemos
{\footnotesize
\begin{align}
& \nonumber \mathfrak{B}_{n}^{[m-1,\mu]}(x+y;\lambda)\\
&\nonumber =\displaystyle\sum\limits_{k=0}^{n}(-1)^k
\sum\limits_{j=k}^{n}j!\binom{j+\alpha}{j-k}\binom{n}{j}\frac{\alpha+\beta+2k+1}{(\alpha+\beta+k+1)_{j+1}}    \mathfrak{B}_{n-j}^{[m-1,\mu]}(y;\lambda)P_k^{(\alpha,\beta)}(1-2x)
\end{align}
}
\end{obs}

\vspace{0.25cm}

Al establecer $u=1$ y $v=0$ en el \textit{Teorema \ref{teojac}}, tenemos el siguiente corolario.

\begin{coro}\label{coreul1}
{\footnotesize
\begin{align}
& \nonumber \mathcal{E}_{n}^{[m-1,\mu]}(x+y,b,c;\lambda)\\
&\nonumber =\displaystyle\sum\limits_{k=0}^{n}(-1)^k
\sum\limits_{j=k}^{n}j!(\log c)^{j}\binom{j+\alpha}{j-k}\binom{n}{j}\frac{\alpha+\beta+2k+1}{(\alpha+\beta+k+1)_{j+1}}    \mathcal{E}_{n-j}^{[m-1,\mu]}(y,b,c;\lambda)P_k^{(\alpha,\beta)}(1-2x)
\end{align}
}
\end{coro}

\vspace{0.25cm}

\begin{obs}
Si hacemos $b=c=e$ en el \textit{Corolario \ref{coreul1}} tenemos
{\footnotesize
\begin{align}
& \nonumber \mathfrak{E}_{n}^{[m-1,\mu]}(x+y;\lambda)\\
&\nonumber =\displaystyle\sum\limits_{k=0}^{n}(-1)^k
\sum\limits_{j=k}^{n}j!\binom{j+\alpha}{j-k}\binom{n}{j}\frac{\alpha+\beta+2k+1}{(\alpha+\beta+k+1)_{j+1}}    \mathfrak{E}_{n-j}^{[m-1,\mu]}(y;\lambda)P_k^{(\alpha,\beta)}(1-2x)
\end{align}
} la cual es justamente la Ec.(4.7) de \cite{Ming}.
\end{obs}

\vspace{0.25cm}

Al tomar $u=1$ y $v=1$ en el \textit{Teorema \ref{teojac}}, tenemos el siguiente corolario.

\vspace{0.25cm}

\begin{coro}\label{corgen1}
{\footnotesize
\begin{align}
& \nonumber \mathcal{G}_{n}^{[m-1,\mu]}(x+y,b,c;\lambda)\\
&\nonumber =\displaystyle\sum\limits_{k=0}^{n}(-1)^k
\sum\limits_{j=k}^{n}j!(\log c)^{j}\binom{j+\alpha}{j-k}\binom{n}{j}\frac{\alpha+\beta+2k+1}{(\alpha+\beta+k+1)_{j+1}}    \mathcal{G}_{n-j}^{[m-1,\mu]}(y,b,c;\lambda)P_k^{(\alpha,\beta)}(1-2x)
\end{align}
}
\end{coro}

\vspace{0.25cm}

\begin{obs}
Si hacemos $b=c=e$ en el \textit{Corolario \ref{corgen1}} tenemos
{\footnotesize
\begin{align}
& \nonumber \mathfrak{G}_{n}^{[m-1,\mu]}(x+y;\lambda)\\
&\nonumber =\displaystyle\sum\limits_{k=0}^{n}(-1)^k
\sum\limits_{j=k}^{n}j!\binom{j+\alpha}{j-k}\binom{n}{j}\frac{\alpha+\beta+2k+1}{(\alpha+\beta+k+1)_{j+1}}    \mathfrak{G}_{n-j}^{[m-1,\mu]}(y;\lambda)P_k^{(\alpha,\beta)}(1-2x)
\end{align}
}
\end{obs}

\vspace{0.25cm}

\begin{teo}\label{bergen}
La relaci�n
\begin{equation}\label{eq94}
\mathcal{Q}_{n}^{[m-1,\mu]}(x+y,b,c;\lambda;u,v)=\displaystyle\sum\limits_{k=0}^{n}\sum\limits_{j=k}^{n}\frac{k!(\log c)^{j}}{(k+m)!}\binom{n}{j}\binom{j}{k}\mathcal{Q}_{n-j}^{[m-1,\mu]}(y,b,c;\lambda;u,v)B_{j-k}^{[m-1]}(x)
\end{equation}
se mantiene entre los polinomios $\mathcal{Q}_{n}^{[m-1,\mu]}(x,b,c;\lambda;u,v)$ y los polinomios de Bernoulli generalizados definidos por (\ref{eq5}) (con $\alpha=1$).
\end{teo}

\begin{proof}
Por sustituci�n de (\ref{eq60}) en el lado derecho de (\ref{eq9}), tenemos
\begin{align*}
& \mathcal{Q}_{n}^{[m-1,\mu]}(x+y,b,c;\lambda;u,v)\\
&=\displaystyle\sum \limits_{j=0}^{n} \binom{n}{j}\mathcal{Q}_{j}^{[m-1,\mu]}(y,b,c;\lambda;u,v)(\log c)^{n-j}\sum\limits_{k=0}^{n-j}\binom{n-j}{k}\frac{k!}{(k+m)!}B_{n-j-k}^{[m-1]}(x).\\
\intertext{por (\ref{s2}), tenemos}
& =\displaystyle\sum \limits_{j=0}^{n}\sum\limits_{k=0}^{n-j} \binom{n}{j}\mathcal{Q}_{j}^{[m-1,\mu]}(y,b,c;\lambda;u,v)(\log c)^{n-j}\binom{n-j}{k}\frac{k!}{(k+m)!}B_{n-j-k}^{[m-1]}(x).\\
\intertext{aplicando (\ref{s1}) se sigue que}
& =\displaystyle\sum \limits_{k=0}^{n}\sum\limits_{j=0}^{n-k} \binom{n}{j}\mathcal{Q}_{j}^{[m-1,\mu]}(y,b,c;\lambda;u,v)(\log c)^{n-j}\binom{n-j}{k}\frac{k!}{(k+m)!}B_{n-j-k}^{[m-1]}(x).\\
\intertext{aplicando (\ref{s3}), obtenemos}
& =\displaystyle\sum \limits_{k=0}^{n}\sum\limits_{j=k}^{n} \binom{n}{n-j}\mathcal{Q}_{n-j}^{[m-1,\mu]}(y,b,c;\lambda;u,v)(\log c)^{j}\binom{j}{k}\frac{k!}{(k+m)!}B_{j-k}^{[m-1]}(x).\\
\intertext{Finalmente, usando (\ref{comb}) la prueba se completa}
& =\displaystyle\sum\limits_{k=0}^{n}\sum\limits_{j=k}^{n}\frac{k!(\log c)^{j}}{(k+m)!}\binom{n}{j}\binom{j}{k}\mathcal{Q}_{n-j}^{[m-1,\mu]}(y,b,c;\lambda;u,v)B_{j-k}^{[m-1]}(x)
\end{align*}
\end{proof}

\vspace{0.25cm}

Estableciendo $\lambda=:-\lambda, u=0$ y $v=1$ en el \textit{Teorema \ref{bergen}}, y luego multiplicando por $(-1)^\mu$ tenemos el siguiente corolario.

\vspace{0.25cm}

\begin{coro}\label{corber2}
\begin{equation}\label{eq934}
\mathcal{B}_{n}^{[m-1,\mu]}(x+y,b,c;\lambda)=\displaystyle\sum\limits_{k=0}^{n}\sum\limits_{j=k}^{n}\frac{k!(\log c)^{j}}{(k+m)!}\binom{n}{j}\binom{j}{k}\mathcal{B}_{n-j}^{[m-1,\mu]}(y,b,c;\lambda)B_{j-k}^{[m-1]}(x)
\end{equation}
\end{coro}

\vspace{0.25cm}

\begin{obs}
Si hacemos $b=c=e$ en el \textit{Corolario \ref{corber2}} tenemos
\begin{equation}
\mathfrak{B}_{n}^{[m-1,\mu]}(x+y;\lambda)=\displaystyle\sum\limits_{k=0}^{n}\sum\limits_{j=k}^{n}\frac{k!}{(k+m)!}\binom{n}{j}\binom{j}{k}\mathfrak{B}_{n-j}^{[m-1,\mu]}(y;\lambda)B_{j-k}^{[m-1]}(x)
\end{equation}
\end{obs}

\vspace{0.25cm}

Estableciendo $u=1$ y $v=0$ en el \textit{Teorema \ref{bergen}}, tenemos el siguiente corolario.

\vspace{0.25cm}

\begin{coro}\label{coreul2}
\begin{equation}
\mathcal{E}_{n}^{[m-1,\mu]}(x+y,b,c;\lambda)=\displaystyle\sum\limits_{k=0}^{n}\sum\limits_{j=k}^{n}\frac{k!(\log c)^{j}}{(k+m)!}\binom{n}{j}\binom{j}{k}\mathcal{E}_{n-j}^{[m-1,\mu]}(y,b,c;\lambda)B_{j-k}^{[m-1]}(x)
\end{equation}
\end{coro}

\vspace{0.25cm}

\begin{obs}
Si hacemos $b=c=e$ en el \textit{Corolario \ref{coreul2}} tenemos
\begin{equation}
\mathfrak{E}_{n}^{[m-1,\mu]}(x+y;\lambda)=\displaystyle\sum\limits_{k=0}^{n}\sum\limits_{j=k}^{n}\frac{k!}{(k+m)!}\binom{n}{j}\binom{j}{k}\mathfrak{E}_{n-j}^{[m-1,\mu]}(y;\lambda)B_{j-k}^{[m-1]}(x)
\end{equation}
la cual es la Ec.(4.11) de \cite{Ming}.
\end{obs}

\vspace{0.25cm}

Establceciendo $u=1$ y $v=1$ en el \textit{Teorema \ref{bergen}}, tenemos el siguiente corolario.

\vspace{0.25cm}

\begin{coro}\label{corgen2}
\begin{equation}
\mathcal{G}_{n}^{[m-1,\mu]}(x+y,b,c;\lambda)=\displaystyle\sum\limits_{k=0}^{n}\sum\limits_{j=k}^{n}\frac{k!(\log c)^{j}}{(k+m)!}\binom{n}{j}\binom{j}{k}\mathcal{G}_{n-j}^{[m-1,\mu]}(y,b,c;\lambda)B_{j-k}^{[m-1]}(x)
\end{equation}
\end{coro}

\vspace{0.25cm}

\begin{obs}
Si hacemos $b=c=e$ en el \textit{Corolario \ref{corgen2}} tenemos
\begin{equation}
\mathfrak{G}_{n}^{[m-1,\mu]}(x+y;\lambda)=\displaystyle\sum\limits_{k=0}^{n}\sum\limits_{j=k}^{n}\frac{k!}{(k+m)!}\binom{n}{j}\binom{j}{k}\mathfrak{G}_{n-j}^{[m-1,\mu]}(y;\lambda)B_{j-k}^{[m-1]}(x)
\end{equation}
\end{obs}

\vspace{0.25cm}

\begin{teo}\label{teostir}
La relaci�n
\begin{equation}\label{eq14}
\mathcal{Q}_{n}^{[m-1,\mu]}(x+y,b,c;\lambda;u,v)=\displaystyle\sum\limits_{k=0}^{n}k!
\binom{x}{k}\sum\limits_{j=k}^{n}\binom{n}{j}\mathcal{Q}_{n-j}^{[m-1,\mu]}(y,b,c;\lambda;u,v)(\log c)^{j}\mathfrak{S}(j,k)
\end{equation}
se mantiene entre los polinomios $\mathcal{Q}_{n}^{[m-1,\mu]}(x,b,c;\lambda;u,v)$ y los n�meros de Stirling de segunda clase $\mathfrak{S}(n,k)$ definidos en (\ref{eq15}).
\end{teo}

\begin{proof}
Por sustituci�n de (\ref{eq15}) en el lado derecho de (\ref{eq9}), tenemos
\begin{align*}
& \mathcal{Q}_{n}^{[m-1,\mu]}(x+y,b,c;\lambda;u,v)=\displaystyle\sum \limits_{j=0}^{n} \binom{n}{j}\mathcal{Q}_{j}^{[m-1,\mu]}(y,b,c;\lambda;u,v)(\log c)^{n-j}\sum\limits_{k=0}^{n-j}\binom{x}{k}k!\mathfrak{S}(n-j,k)\\
\intertext{por (\ref{s2}) se sigue que}
& =\displaystyle\sum \limits_{j=0}^{n}\sum\limits_{k=0}^{n-j} \binom{n}{j}\mathcal{Q}_{j}^{[m-1,\mu]}(y,b,c;\lambda;u,v)(\log c)^{n-j}\binom{x}{k}k!\mathfrak{S}(n-j,k)\\
\intertext{ahora, por (\ref{s1}) entonces}
& =\displaystyle\sum \limits_{k=0}^{n}\sum\limits_{j=0}^{n-k} \binom{n}{j}\mathcal{Q}_{j}^{[m-1,\mu]}(y,b,c;\lambda;u,v)(\log c)^{n-j}\binom{x}{k}k!\mathfrak{S}(n-j,k)\\
\intertext{usando (\ref{s2}) nuevamente, tenemos}
& =\displaystyle\sum \limits_{k=0}^{n}\binom{x}{k}k!\sum\limits_{j=0}^{n-k} \binom{n}{j}\mathcal{Q}_{j}^{[m-1,\mu]}(y,b,c;\lambda;u,v)(\log c)^{n-j}\mathfrak{S}(n-j,k)\\
\intertext{usando (\ref{s3}) , obtenemos}
& =\displaystyle\sum \limits_{k=0}^{n}\binom{x}{k}k!\sum\limits_{j=k}^{n} \binom{n}{n-j}\mathcal{Q}_{n-j}^{[m-1,\mu]}(y,b,c;\lambda;u,v)(\log c)^{j}\mathfrak{S}(j,k).\\
\intertext{Finalmente, por (\ref{comb}) se completa la prueba}
& =\displaystyle\sum\limits_{k=0}^{n}k!
\binom{x}{k}\sum\limits_{j=k}^{n}\binom{n}{j}\mathcal{Q}_{n-j}^{[m-1,\mu]}(y,b,c;\lambda;u,v)(\log c)^{j}\mathfrak{S}(j,k).
\end{align*}
\end{proof}

\vspace{0.25cm}

Estableciendo $\lambda=:-\lambda, u=0$ y $v=1$ en el \textit{Teorema \ref{teostir}}, y luego multiplicando por $(-1)^\mu$ tenemos el siguiente corolario.

\vspace{0.25cm}

\begin{coro}\label{corst2}
\begin{equation}\label{st934}
\mathcal{B}_{n}^{[m-1,\mu]}(x+y,b,c;\lambda)=\displaystyle\sum\limits_{k=0}^{n}k!
\binom{x}{k}\sum\limits_{j=k}^{n}\binom{n}{j}\mathcal{B}_{n-j}^{[m-1,\mu]}(y,b,c;\lambda)(\log c)^{j}\mathfrak{S}(j,k)
\end{equation}
\end{coro}

\vspace{0.25cm}

\begin{obs}
Si hacemos $b=c=e$ en el \textit{Corolario \ref{corst2}} tenemos
\begin{equation}
\mathfrak{B}_{n}^{[m-1,\mu]}(x+y;\lambda)=\displaystyle\sum\limits_{k=0}^{n}k!
\binom{x}{k}\sum\limits_{j=k}^{n}\binom{n}{j}\mathfrak{B}_{n-j}^{[m-1,\mu]}(y;\lambda)\mathfrak{S}(j,k)
\end{equation}
\end{obs}

\vspace{0.25cm}

Al establecer $u=1$ y $v=0$ en el \textit{Teorema \ref{teostir}}, tenemos el siguiente corolario.

\vspace{0.25cm}

\begin{coro}\label{corstl2}
\begin{equation}
\mathcal{E}_{n}^{[m-1,\mu]}(x+y,b,c;\lambda)=\displaystyle\sum\limits_{k=0}^{n}k!
\binom{x}{k}\sum\limits_{j=k}^{n}\binom{n}{j}\mathcal{E}_{n-j}^{[m-1,\mu]}(y,b,c;\lambda)(\log c)^{j}\mathfrak{S}(j,k)
\end{equation}
\end{coro}

\vspace{0.25cm}

\begin{obs}
Si hacemos $b=c=e$ en el \textit{Corolario \ref{corstl2}} tenemos
\begin{equation}
\mathfrak{E}_{n}^{[m-1,\mu]}(x+y;\lambda)=\displaystyle\sum\limits_{k=0}^{n}k!
\binom{x}{k}\sum\limits_{j=k}^{n}\binom{n}{j}\mathfrak{E}_{n-j}^{[m-1,\mu]}(y;\lambda)\mathfrak{S}(j,k)
\end{equation}
la cual es justamente la Ec.(4.3) de \cite{Ming}.
\end{obs}

\vspace{0.25cm}

Haciendo $u=1$ y $v=1$ en el \textit{Teorema \ref{teostir}}, tenemos el siguiente corolario.

\vspace{0.25cm}

\begin{coro}\label{corgst2}
\begin{equation}
\mathcal{G}_{n}^{[m-1,\mu]}(x+y,b,c;\lambda)=\displaystyle\sum\limits_{k=0}^{n}k!
\binom{x}{k}\sum\limits_{j=k}^{n}\binom{n}{j}\mathcal{G}_{n-j}^{[m-1,\mu]}(y,b,c;\lambda)(\log c)^{j}\mathfrak{S}(j,k)
\end{equation}
\end{coro}

\vspace{0.25cm}

\begin{obs}
Si hacemos $b=c=e$ en el \textit{Corolario \ref{corgst2}} tenemos
\begin{equation}
\mathfrak{G}_{n}^{[m-1,\mu]}(x+y;\lambda)=\displaystyle\sum\limits_{k=0}^{n}k!
\binom{x}{k}\sum\limits_{j=k}^{n}\binom{n}{j}\mathfrak{G}_{n-j}^{[m-1,\mu]}(y;\lambda)\mathfrak{S}(j,k)
\end{equation}
\end{obs}

\vspace{0.25cm}

\begin{teo}\label{teolag}
La relaci�n
\begin{align}
& \nonumber \mathcal{Q}_{n}^{[m-1,\mu]}(x+y,b,c;\lambda;u,v)\\
& \label{eq20} =\displaystyle
\sum\limits_{k=0}^{n}(-1)^k
\sum\limits_{j=k}^{n}j!\binom{n}{j}\binom{j+\alpha}{j-k}(\log c)^j\mathcal{Q}_{n-j}^{[m-1,\mu]}(y,b,c;\lambda;u,v)L_k^{(\alpha)}(x).
\end{align}
se mantiene entre los polinomios $\mathcal{Q}_{n}^{[m-1,\mu]}(x,b,c;\lambda;u,v)$ y los polinomios de Laguerre $L_n^{(\alpha)}(x)$ definidos por
(\ref{lag}).
\end{teo}

\vspace{0.25cm}

\begin{proof}
Si hacemos $m=n-j$ en (\ref{lg1}), tenemos:
\begin{equation}\label{eq22}
x^{n-j}=(n-j)!\displaystyle\sum\limits_{k=0}^{n-j}(-1)^k
\binom{n-j+\alpha}{n-j-k}L_k^{(\alpha)}(x).
\end{equation}
Ahora, al sustituir (\ref{eq22}) en el miembro de la derecha de (\ref{eq9}) se sigue que:
\begin{align*}
& \mathcal{Q}_{n}^{[m-1,\mu]}(x+y,b,c;\lambda;u,v)\\
& =\displaystyle\sum
\limits_{j=0}^{n}
\binom{n}{j}\mathcal{Q}_{j}^{[m-1,\mu]}(y,b,c;\lambda;u,v)(\log c)^{n-j}(n-j)!\displaystyle\sum\limits_{k=0}^{n-j}(-1)^k
\binom{n-j+\alpha}{n-j-k}L_k^{\alpha}(x)\\
\intertext{por aplicaci\'on de (\ref{s2}), obtenemos}
& =\displaystyle\sum
\limits_{j=0}^{n}\sum\limits_{k=0}^{n-j}
\binom{n}{j}\mathcal{Q}_{j}^{[m-1,\mu]}(y,b,c;\lambda;u,v)(\log c)^{n-j}(n-j)!(-1)^k
\binom{n-j+\alpha}{n-j-k}L_k^{\alpha}(x)\\
\intertext{usando (\ref{s1}), tenemos}
& =\displaystyle\sum
\limits_{k=0}^{n}\sum\limits_{j=0}^{n-k}
\binom{n}{j}\mathcal{Q}_{j}^{[m-1,\mu]}(y,b,c;\lambda;u,v)(\log c)^{n-j}(n-j)!(-1)^k
\binom{n-j+\alpha}{n-j-k}L_k^{\alpha}(x)
\end{align*}
Finalmente al aplicar nuevamente (\ref{s2}) y luego (\ref{s3}) concluimos la prueba.
\end{proof}

\vspace{0.25cm}
Haciendo $x:=k\alpha$, $a=:x$ en (\ref{ch3}) y sustituyendo $L_k^{(\alpha)}(x)$ en el \textit{Teorema \ref{teolag}} tenemos el siguiente corolario.

\vspace{0.25cm}
\begin{coro}
La relaci�n
\begin{align}
& \nonumber \mathcal{Q}_{n}^{[m-1,\mu]}(x+y,b,c;\lambda;u,v)\\
& \label{eq201} =\displaystyle
\sum\limits_{k=0}^{n}\frac{(-1)^k}{k!}
\sum\limits_{j=k}^{n}j!\binom{n}{j}\binom{j+\alpha}{j-k}(\log c)^j\mathcal{Q}_{n-j}^{[m-1,\mu]}(y,b,c;\lambda;u,v)C_k^{(x)}(k+\alpha).
\end{align}
se mantiene entre los polinomios $\mathcal{Q}_{n}^{[m-1,\mu]}(x,b,c;\lambda;u,v)$ y los polinomios de Charlier $C_n^{(a)}(x)$ definidos por
(\ref{chr1}).
\end{coro}


\vspace{0.25cm}

Por otro lado estableciendo $\lambda=:-\lambda, u=0$ y $v=1$ en el \textit{Teorema \ref{teolag}}, y luego multiplicando por $(-1)^\mu$ tenemos el siguiente corolario.

\vspace{0.25cm}

\begin{coro}\label{corst2l}
\begin{align}
& \nonumber \mathcal{B}_{n}^{[m-1,\mu]}(x+y,b,c;\lambda)\\
& \label{st934l} =\displaystyle
\sum\limits_{k=0}^{n}(-1)^k
\sum\limits_{j=k}^{n}j!\binom{n}{j}\binom{j+\alpha}{j-k}(\log c)^j\mathcal{B}_{n-j}^{[m-1,\mu]}(y,b,c;\lambda)L_k^{(\alpha)}(x).
\end{align}
\end{coro}

\vspace{0.25cm}

\begin{obs}
Si hacemos $b=c=e$ en el \textit{Corolario \ref{corst2l}} tenemos
\begin{align}
& \nonumber \mathfrak{B}_{n}^{[m-1,\mu]}(x+y;\lambda)\\
& \label{uil90} =\displaystyle
\sum\limits_{k=0}^{n}(-1)^k
\sum\limits_{j=k}^{n}j!\binom{n}{j}\binom{j+\alpha}{j-k}\mathfrak{B}_{n-j}^{[m-1,\mu]}(y;\lambda)L_k^{(\alpha)}(x).
\end{align}
\end{obs}

\vspace{0.25cm}

Al establecer $u=1$ y $v=0$ en el \textit{Teorema \ref{teolag}}, tenemos el siguiente corolario.

\vspace{0.25cm}

\begin{coro}
\begin{align}
& \nonumber \mathcal{E}_{n}^{[m-1,\mu]}(x+y,b,c;\lambda)\\
& \label{corstl2l} =\displaystyle
\sum\limits_{k=0}^{n}(-1)^k
\sum\limits_{j=k}^{n}j!\binom{n}{j}\binom{j+\alpha}{j-k}(\log c)^j\mathcal{E}_{n-j}^{[m-1,\mu]}(y,b,c;\lambda)L_k^{(\alpha)}(x).
\end{align}
\end{coro}

\vspace{0.25cm}

\begin{obs}
Si hacemos $b=c=e$ en el \textit{Corolario \ref{corstl2l}} tenemos
\begin{align}
& \nonumber \mathfrak{E}_{n}^{[m-1,\mu]}(x+y;\lambda)\\
& \label{uil901} =\displaystyle
\sum\limits_{k=0}^{n}(-1)^k
\sum\limits_{j=k}^{n}j!\binom{n}{j}\binom{j+\alpha}{j-k}\mathfrak{E}_{n-j}^{[m-1,\mu]}(y;\lambda)L_k^{(\alpha)}(x).
\end{align}
la cual es justamente la Ec.(4.5) de \cite{Ming}.
\end{obs}

\vspace{0.25cm}

Haciendo $u=1$ y $v=1$ en el \textit{Teorema \ref{teolag}}, tenemos el siguiente corolario.

\vspace{0.25cm}

\begin{coro}
\begin{align}
& \nonumber \mathcal{G}_{n}^{[m-1,\mu]}(x+y,b,c;\lambda)\\
& \label{corgst2l} =\displaystyle
\sum\limits_{k=0}^{n}(-1)^k
\sum\limits_{j=k}^{n}j!\binom{n}{j}\binom{j+\alpha}{j-k}(\log c)^j\mathcal{G}_{n-j}^{[m-1,\mu]}(y,b,c;\lambda)L_k^{(\alpha)}(x).
\end{align}
\end{coro}

\vspace{0.25cm}

\begin{obs}
Si hacemos $b=c=e$ en el \textit{Corolario \ref{corgst2l}} tenemos
\begin{align}
& \nonumber \mathfrak{G}_{n}^{[m-1,\mu]}(x+y;\lambda)\\
& \label{uil902} =\displaystyle
\sum\limits_{k=0}^{n}(-1)^k
\sum\limits_{j=k}^{n}j!\binom{n}{j}\binom{j+\alpha}{j-k}\mathfrak{G}_{n-j}^{[m-1,\mu]}(y;\lambda)L_k^{(\alpha)}(x).
\end{align}
\end{obs}

\vspace{0.25cm}

\begin{teo}\label{teoherm}
La relaci�n
\begin{align}
& \nonumber \mathcal{Q}_{n}^{[m-1,\mu]}(x+y,b,c;\lambda;u,v)
\\
& \label{eq23} =\displaystyle\sum\limits_{j=0}^{n}\sum\limits_{k=0}^{[(n-j)/2]}2^{-(n-j)}\binom{n}{j}\binom{n-j}{2k}\displaystyle\frac{(2k)!}{k!}(\log c)^{n-j}\mathcal{Q}_{j}^{[m-1,\mu]}(y,b,c;\lambda;u,v)H_{n-j-2k}(x).
\end{align}
se mantiene entre los polinomios $\mathcal{Q}_{n}^{[m-1,\mu]}(x,b,c;\lambda;u,v)$ y los polinomios de Hermite $H_n(x)$ definidos por
(\ref{her}).
\end{teo}

\vspace{0.25cm}

\begin{proof}
Consideramos la expresi\'on (\ref{her1}), despejamos $x^m$  y hacemos $m=n-j$ tenemos:
\begin{equation}\label{eq25}
x^{n-j}=\displaystyle\sum\limits_{k=0}^{[n-j/2]}2^{-(n-j)}\binom{n-j}{2k}\frac{(2k)!}{k!}H_{n-j-2k}(x),
\end{equation}
finalmente al sustituir (\ref{eq25}) en el miembro de la derecha de (\ref{eq9}) y haciendo uso de (\ref{s2}) finalizamos la prueba.
\end{proof}

\vspace{0.25cm}

Estableciendo $\lambda=:-\lambda, u=0$ y $v=1$ en el \textit{Teorema \ref{teoherm}}, y luego multiplicando por $(-1)^\mu$ tenemos el siguiente corolario.

\vspace{0.25cm}

\begin{coro}\label{huy3}
\begin{align}
& \nonumber \mathcal{B}_{n}^{[m-1,\mu]}(x+y,b,c;\lambda)
\\
& \label{st934h} =\displaystyle\sum\limits_{j=0}^{n}\sum\limits_{k=0}^{[(n-j)/2]}2^{-(n-j)}\binom{n}{j}\binom{n-j}{2k}\displaystyle\frac{(2k)!}{k!}(\log c)^{n-j}\mathcal{B}_{j}^{[m-1,\mu]}(y,b,c;\lambda)H_{n-j-2k}(x).
\end{align}
\end{coro}

\vspace{0.25cm}

\begin{obs}
Si hacemos $b=c=e$ en el \textit{Corolario \ref{huy3}} tenemos
\begin{align}
& \nonumber \mathfrak{B}_{n}^{[m-1,\mu]}(x+y;\lambda)
\\
& \label{st934h1} =\displaystyle\sum\limits_{j=0}^{n}\sum\limits_{k=0}^{[(n-j)/2]}2^{-(n-j)}\binom{n}{j}\binom{n-j}{2k}\displaystyle\frac{(2k)!}{k!}\mathfrak{B}_{j}^{[m-1,\mu]}(y;\lambda)H_{n-j-2k}(x).
\end{align}
\end{obs}

\vspace{0.25cm}

Al establecer $u=1$ y $v=0$ en el \textit{Teorema \ref{teoherm}}, tenemos el siguiente corolario.

\vspace{0.25cm}

\begin{coro}\label{huy2}
\begin{align}
& \nonumber \mathcal{E}_{n}^{[m-1,\mu]}(x+y,b,c;\lambda)
\\
& \label{st934h2} =\displaystyle\sum\limits_{j=0}^{n}\sum\limits_{k=0}^{[(n-j)/2]}2^{-(n-j)}\binom{n}{j}\binom{n-j}{2k}\displaystyle\frac{(2k)!}{k!}(\log c)^{n-j}\mathcal{E}_{j}^{[m-1,\mu]}(y,b,c;\lambda)H_{n-j-2k}(x).
\end{align}
el cual es justamente el \textit{Teorema 4.3} en \cite{Tremblay2}.
\end{coro}

\vspace{0.25cm}

\begin{obs}
Si hacemos $b=c=e$ en el \textit{Corolario \ref{huy2}} tenemos
\begin{align}
& \nonumber \mathfrak{E}_{n}^{[m-1,\mu]}(x+y;\lambda)
\\
& \label{st934h3} =\displaystyle\sum\limits_{j=0}^{n}\sum\limits_{k=0}^{[(n-j)/2]}2^{-(n-j)}\binom{n}{j}\binom{n-j}{2k}\displaystyle\frac{(2k)!}{k!}\mathfrak{E}_{j}^{[m-1,\mu]}(y;\lambda)H_{n-j-2k}(x).
\end{align}
\end{obs}

\vspace{0.25cm}

Haciendo $u=1$ y $v=1$ en el \textit{Teorema \ref{teoherm}}, tenemos el siguiente corolario.

\vspace{0.25cm}

\begin{coro}\label{huy}
\begin{align}
& \nonumber \mathcal{G}_{n}^{[m-1,\mu]}(x+y,b,c;\lambda)
\\
& \label{st934h4} =\displaystyle\sum\limits_{j=0}^{n}\sum\limits_{k=0}^{[(n-j)/2]}2^{-(n-j)}\binom{n}{j}\binom{n-j}{2k}\displaystyle\frac{(2k)!}{k!}(\log c)^{n-j}\mathcal{G}_{j}^{[m-1,\mu]}(y,b,c;\lambda)H_{n-j-2k}(x).
\end{align}
\end{coro}

\vspace{0.25cm}

\begin{obs}
Si hacemos $b=c=e$ en el \textit{Corolario \ref{huy}} tenemos
\begin{align}
& \nonumber \mathfrak{G}_{n}^{[m-1,\mu]}(x+y;\lambda)
\\
& \label{st934h5} =\displaystyle\sum\limits_{j=0}^{n}\sum\limits_{k=0}^{[(n-j)/2]}2^{-(n-j)}\binom{n}{j}\binom{n-j}{2k}\displaystyle\frac{(2k)!}{k!}\mathfrak{G}_{j}^{[m-1,\mu]}(y;\lambda)H_{n-j-2k}(x).
\end{align}
\end{obs}

\vspace{0.25cm}

\begin{teo}\label{teobess}
La relaci�n
\begin{align}
& \nonumber \mathcal{Q}_{n}^{[m-1,\mu]}(x+y,b,d;\lambda;u,v)
\\
& \label{eq23b} =\displaystyle\sum\limits_{k=0}^{n}(-1)^kk!(c+2k)\varphi_k (x,c)\sum\limits_{j=k}^{n}\binom{n}{j}\binom{j}{k}\displaystyle\frac{(\log d)^{j}\mathcal{Q}_{n-j}^{[m-1,\mu]}(y,b,d;\lambda;u,v)}{(c)_{j+k+1}}.
\end{align}
se mantiene entre los polinomios $\mathcal{Q}_{n}^{[m-1,\mu]}(x,b,d;\lambda;u,v)$ y los polinomios de Bessel generalizados $\varphi_n (x,c)$ definidos por
(\ref{sta}).
\end{teo}

\vspace{0.25cm}

\begin{proof}
De acuerdo a la ecuaci�n (\ref{lbess}), haciendo $m=n-j$, tenemos:
\begin{align}
& \label{bess1} x^{n-j}=\displaystyle \sum\limits_{k=0}^{n-j}\binom{n-j}{k}k!\frac{(-1)^k(c+2k)\varphi_k (x,c)}{(c)_{n-j+k+1}}\\
\intertext{sustituyendo (\ref{bess1}) en el miembro de la derecha de (\ref{eq9}), tenemos}
& \nonumber \mathcal{Q}_{n}^{[m-1,\mu]}(x+y,b,d;\lambda;u,v)
\\
& \nonumber =\displaystyle\sum\limits_{j=0}^{n}\binom{n}{j}\mathcal{Q}_{j}^{[m-1,\mu]}(y,b,d;\lambda;u,v)(\log d)^{n-j}\sum\limits_{k=0}^{n-j}\binom{n-j}{k}k!\frac{(-1)^k(c+2k)\varphi_k (x,c)}{(c)_{n-j+k+1}}\\
\intertext{aplicando (\ref{s2}), tenemos}
& \nonumber =\displaystyle\sum\limits_{j=0}^{n}\sum\limits_{k=0}^{n-j}\binom{n}{j}\mathcal{Q}_{j}^{[m-1,\mu]}(y,b,d;\lambda;u,v)(\log d)^{n-j}\binom{n-j}{k}k!\frac{(-1)^k(c+2k)\varphi_k (x,c)}{(c)_{n-j+k+1}}\\
\intertext{usando (\ref{s1}), tenemos}
& \nonumber =\displaystyle\sum\limits_{k=0}^{n}\sum\limits_{j=0}^{n-k}\binom{n}{j}\mathcal{Q}_{j}^{[m-1,\mu]}(y,b,d;\lambda;u,v)(\log d)^{n-j}\binom{n-j}{k}k!\frac{(-1)^k(c+2k)\varphi_k (x,c)}{(c)_{n-j+k+1}}\\
\intertext{por (\ref{s3}), obtenemos}
& \nonumber =\displaystyle\sum\limits_{k=0}^{n}\sum\limits_{j=k}^{n}\binom{n}{n-j}\mathcal{Q}_{n-j}^{[m-1,\mu]}(y,b,d;\lambda;u,v)(\log d)^{j}\binom{j}{k}k!\frac{(-1)^k(c+2k)\varphi_k (x,c)}{(c)_{j+k+1}}.
\end{align}
Finalmente usando (\ref{s2}) nuevamente y (\ref{comb}) se concluye la prueba.
\end{proof}

\vspace{0.25cm}

Haciendo $c=a-1$, $x:=(-x/b_0)$ y amplificando por $\displaystyle \frac{k!}{(a-1)_k}$ en el \textit{Teorema \ref{teobess}}, tenemos el siguiente corolario.

\vspace{0.25cm}

\begin{coro}
La relaci�n
\begin{align}
& \nonumber \mathcal{Q}_{n}^{[m-1,\mu]}(-x/b_0+y,b,d;\lambda;u,v)=\displaystyle\sum\limits_{k=0}^{n}(-1)^k\Gamma(a-1+k)(a-1+2k)y_k(a,b_0,x)\\
& \label{eq23b1}\times\sum\limits_{j=k}^{n}\binom{n}{j}\binom{j}{k}(\log d)^{j}\frac{\mathcal{Q}_{n-j}^{[m-1,\mu]}(y,b,d;\lambda;u,v)}{\Gamma(a+j+k)}.
\end{align}

se mantiene entre los polinomios $\mathcal{Q}_{n}^{[m-1,\mu]}(x,b,d;\lambda;u,v)$ y los polinomios de Bessel generalizados $y_n(a,b_0,x)$ que son soluci�n de (\ref{difeq1}).
\end{coro}

\vspace{0.25cm}

Haciendo $c=1$, $x:=(-x/2)$ y como $(1)_n=n!$ en el \textit{Teorema \ref{teobess}}, tenemos el siguiente corolario.

\vspace{0.25cm}

\begin{coro}
La relaci�n
\begin{align}
& \nonumber \mathcal{Q}_{n}^{[m-1,\mu]}(-x/2+y,b,d;\lambda;u,v)=\displaystyle\sum\limits_{k=0}^{n}(-1)^kk!(2k+1)y_k(x)\\
& \label{eq23b2}\times\sum\limits_{j=k}^{n}\binom{n}{j}\binom{j}{k}(\log d)^{j}\frac{\mathcal{Q}_{n-j}^{[m-1,\mu]}(y,b,d;\lambda;u,v)}{(j+k+1)!}.
\end{align}

se mantiene entre los polinomios $\mathcal{Q}_{n}^{[m-1,\mu]}(x,b,d;\lambda;u,v)$ y los polinomios de Bessel $y_n(x)$ que son soluci�n de (\ref{difeq}).
\end{coro}

\vspace{0.25cm}


\begin{obs}
Si establecemos $m=1$ y $b=c=e$ en (\ref{eq11}), (\ref{eq17}), (\ref{eq94}), (\ref{eq14}),(\ref{eq20}), (\ref{eq201}), (\ref{eq23}) y  $m=1$, $b=d=e$ en (\ref{eq23b}) entonces se obtienen los resultados correspondientes de los polinomios tipo Apostol generalizados $\mathcal{F}_{n}^{(\mu)}(x;\lambda;u,v)$.
\end{obs}

\vspace{0.25cm}

\section{La nueva clase de polinomios tipo Apostol generalizados basados en Hermite}\label{3g}
Finalmente, en esta secci�n damos una generalizaci�n de los polinomios tipo Apostol generalizados.\\

Consideremos la identidad (\ref{hermite5}), si hacemos $a=0$ y $b=t\log c$, tenemos:
\begin{equation}\label{h6}
\exp\left(y\frac{\partial^2}{\partial x^2}\right)\{c^{xt}\}=c^{xt+(\log c)yt^2}.
\end{equation}
Aplicando el operador $\exp\left(y\frac{\partial^2}{\partial x^2}\right)$ en (\ref{newpol}) y  por la identidad (\ref{h6}) podemos definir la nueva clase de polinomios tipo Apostol generalizados basados en Hermite ${_H}\mathcal{Q}_{n}^{[m-1,\alpha]}(x,y,b,c;\lambda;u,v)$ por la siguiente funci�n generatriz
\begin{equation}\label{new787}
\displaystyle\left(\frac{(2^ut^v)^{m}}{\lambda b^t+\sum\limits_{l=0}^{m-1}\frac{(t\log b)^l}{l!}}\right)^{\alpha}c^{xt+(\log c)yt^2} =\displaystyle\sum\limits_{n=0}^{\infty}
{_H}\mathcal{Q}_{n}^{[m-1,\alpha]}(x,y,b,c;\lambda;u,v)\frac{t^n}{n!}
\end{equation}
con $|t\log b|<|log(-\lambda)|$.

\vspace{0.25cm}

\begin{obs}
Si hacemos $m=1$ y $b=c=e$ en (\ref{new787}) tenemos ${_H}\mathcal{Q}_{n}^{(\alpha)}(x,y,e,e;\lambda;u,v)={_H}\mathcal{F}_{n}^{(\alpha)}(x,y;\lambda;u,v)$ la cual es justamente la Ec.(4.5) de \cite{atp1}.
\end{obs}

\vspace{0.25cm}

\begin{teo}
La nueva clase de polinomios tipo Apostol generalizados basados en Hermite ${_{H}}\mathcal{Q}_n^{[m-1,\alpha]}(x,b,c;\lambda;u,v)$ satisfacen las siguientes relaciones:
\begin{align}
\label{eq823i}
{_{H}}\mathcal{Q}_n^{[m-1,\alpha]}(x,b,c;\lambda;u,v)=& \displaystyle\sum\limits_{j=0}^{n}
\binom{n}{j}\mathcal{Q}_{n-j}^{[m-1,\alpha-1]}(b,c;\lambda;u,v){_{H}}\mathcal{Q}_j^{[m-1]}(x,b,c;\lambda;u,v),\\
\label{eq349i}
{_{H}}\mathcal{Q}_n^{[m-1,\alpha]}(x,b,c;\lambda;u,v)=& \displaystyle\sum\limits_{j=0}^{n}
\binom{n}{j}\mathcal{Q}_{n-j}^{[m-1,\alpha]}(b,c;\lambda;u,v)(\log c)^{j}H_n(x,y),
\end{align}
\end{teo}

\begin{proof}
Aplicando el operador $\exp\left(y\frac{\partial^2}{\partial x^2}\right)$  en (\ref{eq823}) y (\ref{eq349}) y luego usando (\ref{hermite4}) en el miembro de la derecha de (\ref{eq349}) la prueba se concluye.
\end{proof}

\vspace{0.25cm}

\begin{obs}
Si hacemos $m=1$ y $b=c=e$ en (\ref{eq823i}) y (\ref{eq349i}) tenemos  
\begin{align}
\label{eq349ij}
{_{H}}\mathcal{F}_n^{(\alpha)}(x;\lambda;u,v)=& \displaystyle\sum\limits_{j=0}^{n}
\binom{n}{j}\mathcal{F}_{n-j}^{(\alpha)}(\lambda;u,v)H_n(x,y),\\
\label{eq823ij}
{_{H}}\mathcal{F}_n^{(\alpha)}(x;\lambda;u,v)=& \displaystyle\sum\limits_{j=0}^{n}
\binom{n}{j}\mathcal{F}_{n-j}^{(\alpha-1)}(b,c;\lambda;u,v){_{H}}\mathcal{F}_j^{[m-1]}(x;\lambda;u,v),
\end{align}
las cuales son justamente las Ecuaciones (4.9) y (4.10) de \cite{atp1} respectivamente.
\end{obs}

\vspace{0.25cm}

Podemos definir las nuevas clases de los polinomios de Bernoulli basados en Hermite, Euler basados en Hermite y Genocchi basados en Hermite as�:
\begin{align}
\label{newji}
\displaystyle\left(\frac{t^{m}}{\lambda b^t-\sum\limits_{l=0}^{m-1}\frac{(t\log b)^l}{l!}}\right)^{\alpha}c^{xt+(\log c)yt^2} = & \displaystyle\sum\limits_{n=0}^{\infty}
{_H}\mathcal{B}_{n}^{[m-1,\alpha]}(x,b,c;\lambda)\frac{t^n}{n!}\\
\intertext{con $|t\log b|<|log(\lambda)|$,}
\label{newji1} \displaystyle\left(\frac{2^{m}}{\lambda b^t+\sum\limits_{l=0}^{m-1}\frac{(t\log b)^l}{l!}}\right)^{\alpha}c^{xt+(\log c)yt^2} = & \displaystyle\sum\limits_{n=0}^{\infty}
{_H}\mathcal{E}_{n}^{[m-1,\alpha]}(x,b,c;\lambda)\frac{t^n}{n!}\\
\intertext{con $|t\log b|<|log(-\lambda)|$,}
\label{newji2}\displaystyle\left(\frac{2^{m}t^{m}}{\lambda b^t+\sum\limits_{l=0}^{m-1}\frac{(t\log b)^l}{l!}}\right)^{\alpha}c^{xt+(\log c)yt^2} = & \displaystyle\sum\limits_{n=0}^{\infty}
{_H}\mathcal{G}_{n}^{[m-1,\alpha]}(x,b,c;\lambda)\frac{t^n}{n!}\\
\intertext{con $|t\log b|<|log(-\lambda)|$.}\nonumber
\end{align}




 
\chapter{CONCLUSIONES Y RECOMENDACIONES}
\markboth{CAP\'ITULO 4\quad CONCLUSIONES Y RECOMENDACIONES}{CAP\'ITULO 4\quad CONCLUSIONES Y RECOMENDACIONES}
\label{AB4}

%recuperar numeracion arabica%
\def\thechapter{\arabic{chapter}}

En este \'ultimo cap\'itulo haremos un resumen de las t\'ecnicas y resultados fundamentales de nuestro trabajo y expondremos problemas abiertos y futuras l\'ineas de investigaci\'on relativos a *****.

\section{Conclusiones}
Podemos decir que el resultado m\'as importante de este trabajo es el relacionado con *****, (ver teorema ****).



\section{Problemas Abiertos y Futuras L\'{\i}neas de Investigaci\'on}

De lo expuesto en la secci\'on precedente se desprende que algunos  problemas interesantes que permitir\'{\i}an  continuar con investigaciones relacionadas, ser\'{\i}an:


\textbf{Problema 1.}
Series de Fourier y representaci\'on integral de los polinomios $\mathcal{Q}_n^{[m-1,\alpha]}(x,b,c;\lambda;u,v)$.

\textbf{Problema 2.}
Estudio de las propiedades de los polinomios $\mathcal{Q}_n^{(\alpha)}(x;\lambda;a,b,c;u,v)$ definidos en (\ref{eqnewpol2}) y relaci\'on con otros polinomios y n\'umeros.

\textbf{Problema 3.}
Series de Fourier y representaci\'on integral de los polinomios $\mathcal{Q}_n^{(\alpha)}(x;\lambda;a,b,c;u,v)$ definidos en (\ref{eqnewpol2}).

\textbf{Problema 4.}
Estudio de la nueva clase de polinomios tipo Apostol generalizados basados en Hermite ${_{H}}\mathcal{Q}_n^{[m-1,\alpha]}(x,b,c;\lambda;u,v)$ que est\'an definidos en (\ref{new787}).

\textbf{Problema 5.}
Estudio de otra nueva clase de polinomios tipo Apostol generalizados basados en Hermite ${_{H}}\mathcal{Q}_n^{(\alpha)}(x;\lambda;a,b,c;u,v)$ que est\'an definidos por la siguiente funci\'on generatriz:
\begin{equation}\label{eqnewpol23}
\displaystyle\left(\frac{2^uz^v}{\lambda b^z+a^z}\right)^{\alpha}c^{xz+(\log c)yz^2} =\sum\limits_{n=0}^{\infty}
{_{H}}\mathcal{Q}_n^{(\alpha)}(x;\lambda;a,b,c;u,v)\frac{z^n}{n!} 
\end{equation}




  \addcontentsline{toc}{chapter}{REFERENCIAS BIBLIOGR\'AFICAS} \vspace{1cm}
 
\begin{thebibliography}{10}
\bibitem{AAC} {Atkins and Cheng.} {(1998).} {Optimal regression
designs in the presenceof random block effects.} {JSPI 77, 321-335}.

\bibitem{C} {Cheng.} {Optimal regression designs under random
blocks-efects models.} {(1995).} {Statistica Sinica 5, 485-497}.

\bibitem{DAH} {Debusho and Haines.} {(2007).} {V- and D-optimal
population designs for the simple linear regression model with a random
intercept term.} {Joumal of Statistical Planning and inference 138,
1116-1130}.

\bibitem{EBSAS} {Entholzner, M., Benda, N., Schwabe, R.} {(2005).}
{A note on designs for estimating population parameters.} {Biometrical
Letters 42, 25-41}.

\bibitem{F} {Fedorov, V.V.} {(1972)} {Theory of Optimal Experiments.}
{Academic Press, New York}.

\bibitem{F} {Fedorov, V.V.} {(1997).} {Model-Oriented Design
of Experiments.} {Springer, New York}.

\bibitem{GP} {Goos, Peter.} {(2002)} {The Optimal Design of
Blocked and Split-Plot.} {Springer, New York}.

\bibitem{GAS} {GraBhoff, U and Schwabe, R.} {(2003).} {On the
analysis of paired observations.} {Statistics and Probability Letters
65, 7-12}.

\bibitem{N} {Norell, L.} {(2006).} {Optimal designs for maximun
likelihood estimators in the one-way random model. U.U.D.M.} {Report
2006:24. Department of Mathematics, Uppsala University}.

\bibitem{SAS} {Schmelter, T. and Schwabe, R.} {(2008).} {On
optimal designsin random intercept models.} {Tetra Mountains Mathematical
Publications. 39, 145-53}.

\bibitem{S} {Schwabe, R.} {(1996).} {Optimum Designs for Multi-Factor
Models.} {Springer, New York}.

\bibitem{S} {Silvey, S.D.} {(1980).} {Optimal Designs.} {Chapman
\& Hall, London}.

\bibitem{VCAB} {Van Breukelen, G.J.P., Candel, M.J.J.M. and Berger,
M.P.F.} {(2008).} {Relative efficiency of enequal cluster sizes
for variance component estimation in cluster randomized and multicentre
trials.} {Statiscal Methods in Medical Research, 17, 439-58}.\end{thebibliography}


\end{document}









